\chapter{硬件基础：从需求到机器结构}

先从一个很小的动作开始：把字节 \code{'A'} 交给一个输出设备。

\begin{lstlisting}
output_one_byte('A'):
  wait until the output device can accept one byte
  place 'A' into the device's input slot
\end{lstlisting}

这两行看起来像普通函数调用，其实已经把整台机器压进去了。执行端要知道设备现在能不能接收；设备要记住自己是空闲还是忙；输出动作要有一个边界，不能一半算发生、一半算没发生；如果设备很慢，当前执行流不能永远原地空转；如果输出从 1 字节变成 4096 字节，执行端也不能一直逐字节搬运。

先构造一个最天真的硬件版本：

\begin{lstlisting}
naive_output_A:
  DATA_LINE = 0x41
  VALID_LINE = 1
\end{lstlisting}

它只是把一个电平组合摆到线上。若外部刚好能接收，也许能看到 \code{'A'}；但它没有回答任何关键问题：设备忙时怎么办，什么时候算已经接收，下一次输出 \code{'B'} 时谁改变数据，失败时谁留下错误，复位后这些线应该回到什么状态。它甚至不能区分“这次输出还没开始”和“上一次输出已经结束但线还保持旧值”。

于是第一处缺口出现了：机器不能只有瞬时信号，必须有能保存事实的状态。设备至少要保存 \code{READY}、\code{BUSY}、\code{ERR} 这类状态位；执行端也要保存“现在走到哪一步、手里有哪些值、下一步要做什么”。没有状态，输出不是一个可解释的动作，只是一组随输入变化的电平。

再把硬件想得稍微强一点：

\begin{lstlisting}
polling_output_A:
  while READY == 0:
    keep checking
  DATA = 0x41
\end{lstlisting}

这已经比第一版强，因为它承认设备会忙。但它又暴露出下一批问题。\code{READY} 必须存在于某个可读的状态槽里；\code{DATA} 必须存在于某个可写的数据槽里；\code{while} 需要比较和条件选择；\code{DATA = 0x41} 需要一个明确提交点，不能在设备还没有接收时就让上层相信完成；一直检查 \code{READY} 还会浪费执行资源。也就是说，状态之后马上需要计算、比较、选择下一步和等待事件。

硬件结构就从这些缺口里长出来。先解决“事实会消失”，才有能在时钟边界保存值的单元；再解决“要不要继续、往哪里继续”，才有比较、选择器和下一步位置；再解决“不能永远只输出固定的 \code{'A'}”，才有可被程序控制的取指、译码和写回；再解决“大量数据放不下”和“多个程序不能互相踩”，才有地址、内存和权限；再解决“慢设备和大块数据不能拖住执行端”，才有事件通知和设备侧搬运。

\begin{keyidea}
硬件结构不是零件清单，而是对需求逐层响应的结果。每个硬件名词都要能回答三个问题：它保存了什么事实，允许哪条数据路径发生，失败或完成时给 OS 留下什么证据。
\end{keyidea}

把 \code{output\_one\_byte('A')} 拆成一份需求账本，就能看到硬件能力出现的顺序：

\begin{longtable}{p{0.18\linewidth}p{0.35\linewidth}p{0.35\linewidth}}
\toprule
卡住的地方 & 必须增加的机器能力 & OS 关注原因 \\
\midrule
线上的值会消失 & 状态单元在时钟边界保存上一刻事实 & 启动、异常返回、设备 reset 都要有可恢复状态 \\
不知道设备能不能收 & 设备状态寄存器保存 ready/busy/error & 驱动不能假设设备永远 ready，也不能无限等 \\
不知道下一步选哪条路 & 比较、分支和下一位置选择 & 等待、错误分支、系统调用返回都依赖可恢复控制流 \\
不能执行不同动作 & 指令、取指、译码、寄存器和写回边界 & task 切换和异常处理必须保存软件可见现场 \\
数据放不下 & 能按地址访问一大片存储 & 用户 buffer、栈、堆和文件缓存都不是少量寄存器能装下的 \\
多个程序会互相破坏 & 访问时自动检查地址归属和权限 & 进程隔离、缺页、COW 和用户指针检查必须有硬件根 \\
慢访问拖住执行 & 常用翻译和常用数据保存在更近处，同时保留精确失败边界 & 性能优化不能破坏精确异常和权限边界 \\
设备不是函数 & 外部状态机通过地址窗口、事件和状态位接入机器 & 驱动面对的是异步硬件状态，不是可靠函数返回 \\
大块 I/O 太重 & 用请求记录描述一批数据，让设备侧完成搬运并留下完成证据 & buffer 所有权和设备访问范围必须可验证 \\
\bottomrule
\end{longtable}

若把这条路径简化成“调用一个输出函数”，执行现场、地址、设备状态和异步完成都会失去来源。机器先满足保存、计算、执行、寻址、隔离、加速、通信、通知和恢复这些需求；具体硬件名称只是这些需求被实现后的形状。

每引入一个硬件名词，都按同一套结构读法展开。先看状态保存在哪里；再看数据从哪里进入，经过哪些组合逻辑、队列、互联或内存层次，最后写回哪里；然后看控制信号或权限位怎样选择路径；最后看失败、超时、权限拒绝或完成事件怎样被记录。只有把这四层合在一起，OS 才能判断一个动作是否可重试、能否被抢占、能否被隔离、资源能否释放。

\section{一、从需求增长到现代整机}

组合逻辑只能给出当前输入对应的输出，不能保存“上一刻”。状态寄存器让机器有上一拍和下一拍。寄存器能保存值，但不能决定写入哪个候选结果，于是出现选择器和写使能。能保存和选择后，机器还只能做固定动作；要让不同程序控制不同动作，就需要指令、取指位置、译码和写回。能执行指令后，寄存器容量不够，于是需要内存。内存一出现，多个程序互相破坏的问题就出现，于是需要地址翻译和权限。访问内存太慢，机器把常用翻译和常用数据放到更近处，但仍必须保留可恢复的失败边界。设备一加入，普通内存路径又不够，外部状态机、事件通知和批量搬运路径也必须接进来。

这条链路的重点不在硬件名称增多，而在每增加一项能力，就增加一个 OS 必须管理的边界。能执行指令，就要能保存和恢复现场；能访问地址，就要区分访问权限；能把常用数据留在更近处，就要处理旧副本；能让设备自己搬数据，就要限制设备能访问的内存范围。OS 的复杂性正是这些能力和风险叠加后的结果。

判断任何硬件结构，都要抓住四件事：旧写法卡在哪里，新增结构保存什么状态，数据走哪条路，控制信号怎样选择，失败时留下什么证据。进程、虚拟内存、文件系统和驱动这些 OS 机制，最终都会回到这四件事上。

第一个缺口是稳定状态。纯计算电路只负责把输入组合成输出。若输入稳定，它经过门延迟后给出结果；若输入变化，输出也会跟着变化。这样的组合逻辑没有“上一刻”和“下一刻”的概念。要让机器按步骤运行，必须先有能保存状态的单元在时钟边沿记住结果。

\begin{pdfdiagram}
  \node[os state, minimum width=26mm] (old) at (0,0) {状态寄存器\\old state};
  \node[os compute, minimum width=31mm] (logic) at (38mm,0) {组合逻辑\\计算 next state};
  \node[os state, minimum width=26mm] (new) at (82mm,0) {状态寄存器\\new state};
  \node[os control, minimum width=21mm] (clk) at (38mm,-17mm) {时钟边沿};
  \node[os control, minimum width=21mm] (rst) at (82mm,-17mm) {复位};
  \draw[os strong bus] (old) -- node[os label, above]{当前值} (logic);
  \draw[os strong bus] (logic) -- node[os label, above]{候选下一状态} (new);
  \draw[os bus] (clk) -- node[os label, right]{采样} (logic.south);
  \draw[os feedback] (rst) -- node[os label, right]{拉回已知初值} (new.south);
\end{pdfdiagram}
\diagramnote{机器能按步骤运行，靠的是“状态寄存器保存上一拍，组合逻辑计算下一拍，时钟边沿统一提交”。}

\begin{lstlisting}
one hardware step:
  old_state is held in registers
  combinational logic computes next_state
  on clock edge:
    registers capture next_state
\end{lstlisting}

把它放回开头的输出设备例子，会更容易看清楚。设备的 \code{READY} 不是一句软件判断，而是设备状态寄存器里的一个 bit。假设设备刚复位完成，\code{READY=1}，表示数据寄存器空闲。主控端读取状态地址时，总线把这个已保存的 bit 返回给主控端；主控端写入 \code{'A'} 时，写事务到达设备的数据寄存器。设备在自己的时钟边界接收这个字节，通常会把 \code{READY} 清成 0，表示暂时不能再收下一个字节；等外部发送完成后，设备控制逻辑再把 \code{READY} 置回 1。

\begin{longtable}{p{0.16\linewidth}p{0.36\linewidth}p{0.36\linewidth}}
\toprule
时刻 & 硬件状态变化 & OS 能观察到什么 \\
\midrule
复位后 & 状态寄存器进入 ready、无错误的初值 & 驱动 probe 后能确认设备可用 \\
主控端读状态 & 地址译码选中设备状态寄存器，返回保存的 bit & 轮询代码看到 ready 或 busy \\
主控端写数据 & \code{'A'} 被写入设备数据寄存器 & 写入不是函数返回，而是提交了一个设备命令 \\
设备接收 & 设备状态机清 \code{READY}，开始外部发送 & 下一次轮询可能看到 busy \\
发送完成 & 设备置 \code{READY=1}，必要时产生完成事件 & 等待者可以继续提交下一个字节 \\
发送失败 & 设备置 \code{ERR=1} 或记录错误码 & 驱动必须停止、清错或 reset \\
\bottomrule
\end{longtable}

这张小表说明“状态”不是抽象名词。OS 所谓等待、超时、唤醒、错误恢复，都是围绕这些硬件保存位建立的：读到 \code{READY=0} 不能假装已经完成，读到 \code{ERR=1} 不能继续写数据槽，reset 后不能沿用 reset 前的请求状态。

复位则把一组关键状态放回已知初值。机器上电后从固定入口或固件指定入口开始；设备 reset 后状态位、队列指针和错误位回到手册规定的状态。没有复位边界，OS 无法判断一个保存单元里的值是旧请求、随机上电值，还是有效完成记录。

\begin{longtable}{p{0.2\linewidth}p{0.34\linewidth}p{0.34\linewidth}}
\toprule
硬件时序概念 & 含义 & OS 直接后果 \\
\midrule
时钟边沿 & 保存单元在离散时刻采样输入 & 执行状态可以在指令边界保存和恢复 \\
复位 & 把状态机拉回已知初始状态 & 启动、驱动 probe、设备 reset 都要重建状态 \\
ready/valid & 一方声明数据有效，另一方声明可接收 & 总线、设备队列、completion 都是握手协议 \\
超时 & 预期状态长期不变化 & 驱动不能无限等 ready，必须报错或 reset \\
错误位 & 状态机报告无法继续 & 内核要记录证据、停止新请求、回收 in-flight 状态 \\
\bottomrule
\end{longtable}

因此，硬件基础不能直接从“机器执行程序”开始。先有时钟和复位，才有可解释的状态；先有状态，寄存器和程序计数器才有意义；先有程序计数器，才有连续取指；先有异常和中断入口，OS 才能在硬件出错或设备完成时接管。换句话说，OS 保存的“现场”，本质上就是硬件在这些状态单元里已经离散化、可命名、可恢复的值。

有了“能保存上一拍”的状态单元之后，机器还不能自动变成处理器。一个保存单元只能记住 0 或 1；把很多保存单元并排放在一起，才能形成 8 位、32 位或 64 位寄存器。寄存器只解决“值放在哪里”，没有解决“下一拍写入什么”。下一拍可能保持原值，可能写入加法结果，可能写入从外部读回的数据，也可能在复位时写入固定初值。要在这些候选值之间选择，就需要选择器，常写作 MUX。

\begin{lstlisting}
one controlled register update:
  candidates = {old_value, alu_result, load_data, reset_value}
  selected   = mux(candidates, control_signal)
  on clock edge if write_enable:
    register = selected
\end{lstlisting}

这段伪代码对应的不是软件分支，而是硬件连线：多个候选输入同时接到选择器，控制信号只决定哪一路能到达寄存器输入端。若 \code{write\_enable} 没有打开，时钟到来时寄存器继续保持旧值。由此形成最小数据通路：状态寄存器提供旧值，组合逻辑产生候选结果，选择器决定写回内容，时钟边沿把结果提交成新状态。

举一个具体的寄存器更新例子。假设寄存器 \code{R0} 当前保存 5，另一个寄存器 \code{R1} 保存 3。本周期控制器希望执行 \code{R0 = R0 + R1}。在时钟边沿到来之前，\code{R0} 和 \code{R1} 的输出已经稳定；ALU 的加法通路计算出候选结果 8；MUX 选择 ALU 结果；\code{write\_enable} 打开。真正把 8 写进 \code{R0} 的动作发生在提交边界，而不是 ALU 刚算出 8 的瞬间。

\begin{longtable}{p{0.18\linewidth}p{0.38\linewidth}p{0.32\linewidth}}
\toprule
阶段 & 硬件里发生什么 & 为什么重要 \\
\midrule
旧状态保持 & \code{R0=5, R1=3} 在寄存器输出端稳定 & 组合逻辑有确定输入 \\
候选结果产生 & ALU 看到两个输入，输出 \code{8} & 此时只是候选值，还没改变 \code{R0} \\
控制信号选择 & MUX 选择 ALU 输出，写使能打开 & 决定本周期是否真的更新状态 \\
时钟边沿提交 & \code{R0} 采样输入端，变成 \code{8} & 从这里开始，软件可见状态改变 \\
未提交分支 & 若写使能关闭，\code{R0} 仍是 \code{5} & 异常、停顿、取消都依赖这个边界 \\
\bottomrule
\end{longtable}

所以“提交”是理解 OS 的核心词。上下文切换保存的是已经提交的软件可见寄存器；page fault 要求 faulting 指令还没有提交错误结果；设备 reset 要求旧队列状态不能被误当作新提交。没有这条边界，OS 就无法区分“硬件正在尝试”与“动作已经发生”。

\begin{pdfdiagram}
  \node[os state, minimum width=25mm] (ra) at (0,10mm) {寄存器 A\\当前值};
  \node[os state, minimum width=25mm] (rb) at (0,-10mm) {寄存器 B\\当前值};
  \node[os control, minimum width=25mm] (mux) at (40mm,0) {选择器 MUX\\选输入};
  \node[os compute, minimum width=30mm] (logic) at (82mm,0) {组合逻辑\\加/比/传递};
  \node[os state, minimum width=27mm] (dst) at (126mm,0) {目标寄存器\\下一拍保存};
  \node[os control, minimum width=33mm] (ctrl) at (82mm,-25mm) {控制状态\\op + write enable};
  \draw[os strong bus] (ra.east) -- node[os label, above]{候选值} (mux.west);
  \draw[os strong bus] (rb.east) -- node[os label, below]{候选值} (mux.west);
  \draw[os strong bus] (mux) -- (logic);
  \draw[os strong bus] (logic) -- node[os label, above]{候选结果} (dst);
  \draw[os bus] (ctrl.north) -- node[os label, right]{选择和写使能} (logic.south);
\end{pdfdiagram}
\diagramnote{硬件动作不是“一个函数返回值”，而是状态寄存器、选择器、组合逻辑和写使能共同完成的一次提交。OS 后来依赖的可恢复现场，必须落在这些已经提交的状态上。}

这条数据通路还没有程序，也没有 CPU。它只是说明硬件怎样从“保存一个值”进化到“按控制信号更新一个值”。即使换成 ALU、寄存器文件、程序计数器、地址翻译和设备寄存器，结构也没有变：旧状态经过数据通路形成候选新状态，再由控制信号决定是否提交。不同硬件名词之间真正相连的，是这条状态更新规则。

在最小数据通路里，组合逻辑如果只能“原样传递”，机器最多保存和搬运值；它还不能计算地址偏移，不能判断状态位是否为零，也不能比较两个数。下一项需求就是：给定两个寄存器值，能不能算出结果，能不能比较大小或是否为零。这个需求推出算术逻辑单元，英文常写作 ALU。ALU 接收两个输入和一个操作选择，输出计算结果，并可能更新 zero、carry、sign、overflow 等标志位。\code{add}、\code{sub}、\code{and}、\code{or}、\code{xor}、\code{compare} 等操作，都可以落到 ALU 或相关执行单元上。

理解 ALU 时，不应只把它看成“会算数的黑盒”。它首先是一组执行 bit 运算的门电路。以一位全加器为例，输入是两个 bit 和一个进位，输出是本位结果和下一位进位。

\begin{lstlisting}
full_adder(a, b, carry_in):
  sum_bit = a xor b xor carry_in
  carry_out = (a & b) | (a & carry_in) | (b & carry_in)
  return sum_bit, carry_out
\end{lstlisting}

一位加法器只能算一位。把很多个一位加法器连起来，低位的 \code{carry\_out} 进入高位的 \code{carry\_in}，就能做 32 位或 64 位加法。真实执行单元会用更快的进位预测或并行前缀结构减少等待，但 OS 视角要抓住的是：加法会产生结果，也会产生标志位。减法通常可以转成“加上补码”；比较通常可以转成“做一次减法但不写回结果，只更新标志位”。

\begin{pdfdiagram}
  \node[os state, minimum width=23mm] (ra) at (0,8mm) {寄存器 A\\操作数};
  \node[os state, minimum width=23mm] (rb) at (0,-8mm) {寄存器 B\\操作数};
  \node[os control, minimum width=22mm] (op) at (34mm,23mm) {op 选择\\ADD/SUB/CMP};
  \node[os compute, minimum width=31mm] (alu) at (43mm,0) {ALU\\加法/逻辑/比较};
  \node[os state, minimum width=25mm] (rd) at (86mm,8mm) {目标寄存器\\result};
  \node[os state, minimum width=25mm] (flags) at (86mm,-11mm) {状态寄存器\\flags};
  \draw[os strong bus] (ra.east) -- (alu.west);
  \draw[os strong bus] (rb.east) -- (alu.west);
  \draw[os bus] (op.south) -- node[os label, right]{选择运算} (alu.north);
  \draw[os strong bus] (alu.east) -- node[os label, above]{结果} (rd.west);
  \draw[os bus] (alu.east) -- ++(8mm,0) |- node[os label, below right]{标志位} (flags.west);
\end{pdfdiagram}
\diagramnote{ALU 不单独构成“程序执行”，它只是数据通路中的计算节点；寄存器提供输入并接收结果，状态寄存器保存可被分支和异常路径使用的证据。}

\begin{longtable}{p{0.18\linewidth}p{0.36\linewidth}p{0.34\linewidth}}
\toprule
标志位 & 来源 & 对 OS 和程序的影响 \\
\midrule
zero & 结果所有 bit 都为 0 & 条件分支、循环结束、错误码判断 \\
sign & 结果最高位为 1 & 有符号比较、边界检查 \\
carry & 无符号加法进位或减法借位 & 地址计算、无符号溢出判断、多精度整数 \\
overflow & 有符号结果超出可表示范围 & 调试、运行时检查、语言语义实现 \\
\bottomrule
\end{longtable}

\begin{lstlisting}
alu(a, b, op):
  if op == ADD: result = a + b
  if op == SUB: result = a - b
  if op == AND: result = a & b
  if op == CMP: flags = compare(a, b)
\end{lstlisting}

但能计算还不等于能运行程序。ALU 不知道 \code{a} 和 \code{b} 从哪里来，不知道结果写回哪里，也不知道下一步做什么。它只是在“寄存器 -> 组合逻辑 -> 寄存器”这条通路中提供一种候选结果。要让 ALU 成为可编程机器的一部分，随之出现的需求是：有一组可命名寄存器保存参数和临时值，有选择器把正确寄存器接到 ALU 输入，有写回端口保存结果，有一个状态保存“下一步取哪条指令”。

\begin{longtable}{p{0.2\linewidth}p{0.34\linewidth}p{0.34\linewidth}}
\toprule
部件 & 最小作用 & OS 相关后果 \\
\midrule
ALU & 整数运算、比较、位操作 & 页表权限位、bitmap、核心掩码、条件分支 \\
寄存器文件 & 保存参数、地址和临时值 & 系统调用参数、trap frame、上下文切换 \\
程序计数器 & 指向下一条指令 & 异常返回、信号、调试、调度恢复 \\
状态寄存器 & 保存条件码、模式位、中断位 & 特权切换、关中断、返回用户态 \\
控制器 & 把指令变成控制信号 & 非法指令、特权指令、系统调用入口 \\
\bottomrule
\end{longtable}

所谓“保存执行上下文”，保存的不是一个抽象对象，而是这组软件可见状态：通用寄存器、栈指针、程序计数器、状态寄存器，以及必要的浮点/SIMD 状态。ALU 本身无须保存，它下一周期可以继续计算；必须保存的是用户程序能够观察到的架构状态。

当机器不再只做固定加法，而是要按程序选择动作时，控制信号就不能靠手工拨开关产生。机器需要把一串指令字节解释成“读哪些寄存器、选择哪种 ALU 操作、结果写回哪里、下一条指令在哪里”。这就引出指令、译码、寄存器文件、数据通路和控制信号。以 \code{add rax, rbx} 为例，硬件不是直接“知道加法语义”，而是在一个数据通路里打开一组控制信号。

\begin{lstlisting}
execute add rax, rbx:
  left  = regfile.read(rax)
  right = regfile.read(rbx)
  out   = alu(left, right, ADD)
  regfile.write(rax, out)
  update_flags(out)
\end{lstlisting}

这已经构成执行指令的基本骨架：寄存器文件提供输入，ALU 计算，写回端口保存结果，状态寄存器记录标志位。再加上取指和译码，就能执行不同指令；再加上 load/store 单元，就能访问更多数据；再加上异常入口，就能在执行失败时交给内核。在这个阶段，可以把这套取指、译码、执行和提交状态的硬件整体称为 CPU。

\begin{pdfdiagram}
  \node[os state, minimum width=22mm] (pc) at (0,18mm) {PC/RIP\\下一条指令};
  \node[os memory, minimum width=25mm] (imem) at (35mm,18mm) {指令存储\\取字节};
  \node[os control, minimum width=27mm] (dec) at (72mm,18mm) {译码/控制\\控制信号};
  \node[os control, minimum width=28mm] (nextpc) at (112mm,18mm) {下一 PC\\顺序/分支/异常};
  \node[os state, minimum width=27mm] (regs) at (0,-12mm) {寄存器文件\\读/写端口};
  \node[os compute, minimum width=28mm] (exe) at (37mm,-12mm) {执行单元\\ALU/分支/地址};
  \node[os memory, minimum width=27mm] (lsu) at (75mm,-12mm) {load/store\\数据访问};
  \node[os state, minimum width=25mm] (wb) at (112mm,-12mm) {写回\\寄存器/flags};
  \draw[os strong bus] (pc) -- node[os label, above]{取指} (imem);
  \draw[os strong bus] (imem) -- node[os label, above]{字节} (dec);
  \draw[os bus] (dec) -- node[os label, above]{控制} (nextpc);
  \draw[os bus] (dec.south) -- ++(0,-9mm) -| (regs.north);
  \draw[os strong bus] (regs) -- node[os label, above]{操作数} (exe);
  \draw[os strong bus] (exe) -- node[os label, above]{结果/地址} (lsu);
  \draw[os strong bus] (lsu) -- node[os label, above]{数据/异常} (wb);
  \draw[os bus] (wb.north) -- ++(0,9mm) -| node[os label, near end, right]{提交} (nextpc.south);
  \draw[os feedback] (nextpc.north) -- ++(0,7mm) -| node[os label, near end, left]{更新} (pc.north);
\end{pdfdiagram}
\diagramnote{CPU 不是一个单块黑盒，而是取指、译码、读寄存器、执行、访存、写回和下一 PC 选择共同组成的状态机。OS 能恢复执行，依赖的是这些结构在提交边界留下的架构状态。}

\begin{lstlisting}
control signals for add rax, rbx:
  reg_src0 = rax
  reg_src1 = rbx
  alu_op = ADD
  writeback_dst = rax
  writeback_src = ALU_RESULT
  update_flags = true
  next_rip = rip + instruction_length
\end{lstlisting}

控制信号把“指令语义”落实成硬件动作。MUX 负责按控制信号选择输入：ALU 左输入来自哪个寄存器，右输入来自寄存器还是立即数，写回数据来自 ALU 还是内存，下一 RIP 来自顺序地址、分支目标、异常入口还是中断返回。写回端口和下一 PC 选择是这张图里最重要的提交点：只要结果还没有写回、下一 PC 还没有提交，硬件就能把这条指令当作未完成；一旦写回完成，OS 就必须把它当作已经发生的架构事实。

把 \code{add rax, rbx} 走细一点。假设进入这条指令前，\code{RIP=0x400100}，\code{rax=5}，\code{rbx=3}。取指单元用 \code{RIP} 访问指令字节；译码器识别出这是一条把 \code{rbx} 加到 \code{rax} 的整数指令；寄存器文件打开 \code{rax} 和 \code{rbx} 的读端口；ALU 输出 8；写回端口把 8 写回 \code{rax}；下一 PC 选择器把 \code{RIP} 更新到下一条指令地址。每一步都接在前面的状态上。

\begin{longtable}{p{0.16\linewidth}p{0.36\linewidth}p{0.36\linewidth}}
\toprule
阶段 & 具体硬件动作 & 提交后的可见事实 \\
\midrule
取指 & \code{RIP=0x400100} 送到指令存储或 I-cache & 还没有改变通用寄存器 \\
译码 & 指令字节被解释成读 \code{rax/rbx}、ALU 加法、写 \code{rax} & 控制信号已确定，结果还未出现 \\
读寄存器 & 寄存器文件输出 \code{5} 和 \code{3} & 仍只是读取旧状态 \\
执行 & ALU 计算 \code{5+3=8}，同时更新候选标志位 & 8 还只是候选结果 \\
写回 & 写端口在提交边界把 8 写入 \code{rax} & \code{rax=8} 成为软件可见事实 \\
更新 PC & 下一 PC 选择顺序地址 & 返回用户态或继续执行时从下一条开始 \\
\bottomrule
\end{longtable}

这也是为什么 OS 保存现场时不保存“ALU 里面正在算什么”。若 \code{add} 已提交，现场里看到的是 \code{rax=8} 和下一条 \code{RIP}；若它没有提交，现场里应该仍能回到 \code{rax=5} 和这条指令。硬件内部可以流水化、乱序化，但对 OS 暴露的提交事实必须保持这种清晰边界。

\code{load} 和 \code{store} 比 \code{add} 多一层问题：它们不只使用寄存器和 ALU，还要先算出一个地址，再按这个地址读写外部存储。最小模型中，地址直接指向某个存储位置，暂不考虑隔离和加速。只要读写失败，CPU 就不能把这条指令标记为已经完成。

\begin{lstlisting}
execute load rax, [rbx + 8]:
  effective_addr = regfile.read(rbx) + 8
  if memory_cannot_read(effective_addr):
    raise_load_fault(effective_addr)
  data = memory_read(effective_addr)
  regfile.write(rax, data)
\end{lstlisting}

最后一步的顺序是关键约束：只有地址检查和读数据都成功后，\code{rax} 才能被写回。若先写回一半再报错，内核就无法把这条指令当作“没有发生过”来重试。异常可恢复，依赖硬件在写回前守住提交边界。即使加入地址隔离和缓存，这个提交边界仍然不能被破坏。

仍用一个具体值走一遍。假设 \code{rbx=0x1000}，指令是 \code{load rax, [rbx + 8]}，那么地址生成单元先得到 \code{0x1008}。如果这个地址可读，内存系统返回 8 个字节，写回端口把它们写进 \code{rax}，然后 \code{RIP} 前进。若地址不可读，硬件不能先把随机数据写进 \code{rax}，再告诉内核“出错了”；它必须保存 fault address、错误原因和 faulting RIP，进入异常入口，让内核决定分配页、发信号、杀进程或 panic。

\begin{longtable}{p{0.18\linewidth}p{0.38\linewidth}p{0.32\linewidth}}
\toprule
路径 & 硬件动作 & OS 后续解释 \\
\midrule
load 成功 & 地址生成成功，权限通过，数据返回，\code{rax} 写回 & 指令已经发生，异常处理不应再重试它 \\
地址不存在 & 页表或最小内存模型报告无法读 & 内核可建立映射后重试同一条指令 \\
权限不允许 & 读/写/执行权限检查失败 & 可能是 COW、非法访问或安全违规 \\
设备或物理错误 & 下层返回不可恢复错误或机器检查 & 进入错误处理、隔离或 panic 路径 \\
\bottomrule
\end{longtable}

这张表给出一条通用判断。缺页、COW、用户指针校验、信号恢复、断点调试，本质上都要求硬件能把“失败发生在提交前”说清楚。否则内核不知道应该重试原指令、跳过原指令，还是把进程结束掉。

\begin{longtable}{p{0.18\linewidth}p{0.4\linewidth}p{0.3\linewidth}}
\toprule
已有能力 & 仍然缺少什么 & 后果 \\
\midrule
只有 ALU & 没有寄存器和 PC & 不能保存程序状态，无法连续执行 \\
只有寄存器和 ALU & 没有指令译码 & 不能由程序选择操作，只能固定连线 \\
有译码和 PC & 没有特权检查 & 用户程序可执行敏感动作 \\
有 load/store & 没有访问边界 & 所有程序共享同一片地址，彼此可以破坏 \\
有异常入口 & 没有精确现场 & 内核无法修复 fault 后重试 \\
\bottomrule
\end{longtable}

从“能计算”过渡到 OS，需要补上三条硬件事实：程序能连续执行，因此必须有可恢复现场；程序可能失败，因此必须有异常证据；多个程序需要共存，因此必须有权限和地址隔离。

能按步骤执行之后，下一项需求是提高吞吐。最小的 CPU 模型可以是一条指令在一个较长周期内完成取指、译码、执行、访存和写回。该模型的问题在于周期必须按最慢指令设计，简单加法也要等待 load/store 的最长路径。流水线把一条指令拆成多个阶段，让不同指令重叠运行。

先看不流水化时的低效。假设一条 \code{add} 只需要取指、译码、ALU 和写回，另一条 \code{load} 还需要等待数据 cache。若硬件要求“每条指令从头到尾独占整条路径”，那么 \code{add} 也要按 \code{load} 的最坏延迟留足时间。结果是：取指部件在执行阶段闲着，ALU 在访存阶段闲着，写回端口在取指阶段闲着。机器不是没有硬件，而是同一时刻只允许一条指令使用这些硬件。

\begin{longtable}{p{0.18\linewidth}p{0.38\linewidth}p{0.32\linewidth}}
\toprule
模型 & 一条指令怎样占用硬件 & 问题在哪里 \\
\midrule
单周期模型 & 取指、译码、执行、访存、写回必须在一个长周期内完成 & 周期被最慢路径决定，简单指令也被拖慢 \\
多周期模型 & 一条指令分多拍走完，但仍独占主要路径 & 部件分时复用，吞吐仍低 \\
流水线模型 & 每拍推进一个阶段，多条指令占据不同阶段 & 吞吐提高，但需要处理相关、预测和异常边界 \\
\bottomrule
\end{longtable}

\begin{lstlisting}
classic pipeline:
  IF   fetch instruction at RIP
  ID   decode and read registers
  EX   run ALU or compute address
  MEM  access memory
  WB   write result back
\end{lstlisting}

流水线不是把一条指令切成五个软件函数，而是在硬件阶段之间加入保存中间结果的流水线寄存器。IF/ID 寄存器保存取到的指令字节和下一 PC；ID/EX 寄存器保存译码后的控制信号、源寄存器编号和立即数；EX/MEM 寄存器保存 ALU 结果或有效地址；MEM/WB 寄存器保存准备写回的数据和目标寄存器号。每个时钟边沿，这些阶段寄存器一起向前推进。这样第 1 条指令在 EX 阶段计算时，第 2 条指令可以同时在 ID 阶段读寄存器，第 3 条指令可以同时在 IF 阶段取指。

\begin{longtable}{p{0.15\linewidth}p{0.22\linewidth}p{0.22\linewidth}p{0.25\linewidth}}
\toprule
周期 & 指令 1 & 指令 2 & 关键硬件事实 \\
\midrule
第 1 拍 & IF 取指 & -- & PC 和 I-cache 工作 \\
第 2 拍 & ID 译码/读寄存器 & IF 取指 & IF/ID 寄存器隔开两个阶段 \\
第 3 拍 & EX 执行 & ID 译码/读寄存器 & ALU 和取指前端可并行工作 \\
第 4 拍 & MEM 访存或跳过 & EX 执行 & load/store 才会使用数据 cache \\
第 5 拍 & WB 写回 & MEM 访存或跳过 & 架构状态在写回处提交 \\
\bottomrule
\end{longtable}

这张表解释了为什么流水线一出现，硬件就必须能“停住某些阶段、冲掉某些阶段、转发某些结果”。如果第 2 条指令马上使用第 1 条指令的结果，而第 1 条还没写回，ID 阶段读到的可能是旧寄存器值。硬件可以把 EX 或 MEM 阶段的新结果直接转发给后面的 EX 阶段；如果数据还没回来，就只能插入停顿。流水线提高的是正常路径吞吐，不是取消了依赖关系。

流水线提高吞吐，也制造新问题。这里的 hazard 指“下一步不能按理想节奏继续推进的情况”，不是软件异常。后一条指令可能依赖前一条结果，分支可能猜错，load 可能访问失败，执行单元也可能被前面的长延迟操作占住。

\begin{longtable}{p{0.18\linewidth}p{0.38\linewidth}p{0.32\linewidth}}
\toprule
流水线问题 & 具体例子 & 硬件或 OS 后果 \\
\midrule
数据相关 & \code{add rax, rbx} 后立刻使用 \code{rax} & 转发、停顿或乱序调度 \\
控制相关 & 分支目标还没确定就要取下一条指令 & 分支预测，猜错后冲刷流水线 \\
结构相关 & 多条指令同时需要同一个执行资源 & 排队、发射限制、吞吐下降 \\
异常边界 & load 访问失败时后面指令可能已猜测执行 & 必须丢弃年轻指令并保留精确现场 \\
\bottomrule
\end{longtable}

现代乱序核心进一步把指令拆成微操作，允许先执行不依赖的操作，但最终仍要按程序顺序提交架构状态。这个“按序提交”是 OS 能处理异常的根。

乱序执行解决的是另一个浪费：流水线虽然能重叠阶段，但遇到 cache miss 或长延迟乘除法时，后面那些完全不依赖它的指令也可能被堵住。乱序核心会把已译码的操作放入队列，源操作数准备好、执行单元空闲时就先执行。为了不破坏程序语义，它还要维护一张提交队列，常称为 reorder buffer。每个条目至少记录：这条操作在程序顺序中的位置、目标架构寄存器、计算出的候选结果、是否完成、是否有异常。内部可以先完成年轻指令，但只有队首的老指令满足条件时，才能把结果提交成软件可见状态。

\begin{longtable}{p{0.18\linewidth}p{0.34\linewidth}p{0.36\linewidth}}
\toprule
结构 & 保存什么 & 为什么 OS 依赖它 \\
\midrule
保留站/调度队列 & 已译码但等待操作数或执行单元的微操作 & 允许独立操作绕过慢操作先执行 \\
物理寄存器 & 临时保存未提交或重命名后的结果 & 消除假依赖，但不直接作为 OS 保存对象 \\
提交队列 & 程序顺序、目标寄存器、完成位、异常位 & 保证异常点之前已经提交，异常点之后都可丢弃 \\
store buffer & 已产生但尚未对外可见的写入 & 锁、屏障、MMIO 顺序必须约束它 \\
\bottomrule
\end{longtable}

\begin{lstlisting}
precise fault boundary:
  older instructions have committed
  faulting instruction has not committed its result
  younger instructions are discarded
  saved RIP points to the faulting instruction
  fault address and error code are recorded
\end{lstlisting}

这个边界决定异常是否可恢复。以 \code{load rax, [p]} 访问失败为例，CPU 不能让后面已经猜测执行的写入对用户程序可见；也不能把 RIP 保存到下一条指令，否则内核处理完原因后无法重试原来的 load。虚拟内存中的懒分配、文件映射和 COW 都依赖这种精确、可重试的异常现场。

用一个小程序看清楚这件事：

\begin{lstlisting}
1: add rbx, 1
2: load rax, [p]       // p may fault
3: store [q], rax
\end{lstlisting}

若第 2 条 \code{load} 缺页，第 1 条 \code{add} 可以已经提交，因为它在程序顺序上更老；第 2 条不能把错误数据写入 \code{rax}；第 3 条即使在内部已经猜测执行，也必须被丢弃，不能真的写入 \code{q}。内核拿到的现场应该是：\code{rbx} 已经加一，\code{RIP} 指向第 2 条，fault address 是 \code{p}，错误码说明是 not-present 或权限问题。缺页处理修复页表后返回，CPU 从第 2 条重新执行。这个例子把流水线、乱序和 OS 缺页处理接在了一起。

\begin{longtable}{p{0.2\linewidth}p{0.36\linewidth}p{0.32\linewidth}}
\toprule
机制 & 它带来的性能能力 & OS 依赖或承担的后果 \\
\midrule
流水线 & 多条指令重叠执行 & 异常必须落在可恢复边界 \\
分支预测 & 减少等待分支结果的空泡 & 预测器可能形成侧信道 \\
乱序执行 & 慢访问或长运算期间执行其他独立操作 & 需要精确异常和按序提交 \\
寄存器重命名 & 消除假依赖，提高并行度 & OS 保存架构寄存器，不保存内部物理寄存器 \\
store buffer & 写入先排队，提高吞吐 & 锁、屏障和设备顺序必须显式表达 \\
\bottomrule
\end{longtable}

现代 CPU 对 OS 暴露两层事实：ISA 层承诺寄存器、异常、地址访问和特权规则；微架构层影响性能和安全，包括流水线停顿、分支预测、store buffer、慢访问等待和侧信道。OS 不能忽略微架构影响，否则锁、调度、内存屏障和安全隔离都无法完整解释。

能执行指令之后，程序会遇到容量问题：寄存器数量有限，所有数据不可能都放在执行单元旁边。数组、栈、堆、文件缓存、网络包都需要一大片可寻址存储，这就是内存。最小内存模型可以表示为“地址到字节的表”：

\begin{lstlisting}
memory[address] -> byte
load  rax, [address]  reads bytes from memory
store [address], rax  writes bytes to memory
\end{lstlisting}

这个表面模型背后对应的是一条物理路径。执行单元先算出地址，地址线或片上互联把这个数送到内存子系统；内存控制器根据地址选择通道、bank 和行列；DRAM 返回一整段数据，通常再被切成 CPU 需要的字节或字。对 OS 来说，最初要抓住两件事：地址不是变量名，而是选择某个存储位置的编号；内存读写不是 ALU 内部动作，而是一次可能等待、失败、被缓存或被权限拦住的外部访问。

\begin{longtable}{p{0.18\linewidth}p{0.38\linewidth}p{0.32\linewidth}}
\toprule
从寄存器到内存 & 结构上增加了什么 & 新问题 \\
\midrule
寄存器读写 & 编号很少，位置在 CPU 内部 & 容量不够，不能放大量对象 \\
直接物理内存 & 地址选择 DRAM 中的字节或 cache line & 多程序互相覆盖，装载位置僵硬 \\
带权限的地址翻译 & 虚拟地址先经页表变成物理地址 & 需要维护页表、处理 fault 和旧翻译 \\
带缓存的访问 & 常用物理 cache line 放在核心附近 & 需要处理一致性、伪共享和 DMA 同步 \\
\bottomrule
\end{longtable}

有了内存，程序才能把大对象放到寄存器之外；随之出现的是隔离问题：如果所有程序都直接使用同一片物理地址，一个程序写错地址就能覆盖另一个程序或内核。为了让多个程序共存，机器必须把“程序看到的地址”和“真实内存位置”分开，并在访问时检查权限。这个硬件叫 MMU，memory management unit。它按页表把虚拟地址翻译成物理地址，并检查读、写、执行、用户/内核等权限。

先用最小反例说明为什么直接物理地址不够。假设程序 A 的数组在物理地址 \code{0x2000}，程序 B 的栈也恰好被错误写到 \code{0x2000} 附近。没有翻译和权限时，硬件只看地址数字，不知道“这是 A 的数组”还是“这是 B 的栈”。一次越界 store 会直接改坏另一个程序。OS 当然可以在软件里检查某些系统调用参数，但用户程序普通 load/store 每秒可以执行大量次，不能每条都陷入内核检查。权限必须放到硬件访存路径上，随每一次 load/store 自动执行。

\begin{longtable}{p{0.22\linewidth}p{0.34\linewidth}p{0.32\linewidth}}
\toprule
访问方式 & 硬件看到什么 & 隔离结果 \\
\midrule
直接物理地址 & 指令给出 \code{0x2000}，内存系统直接访问物理位置 & 任意错误地址都可能改坏别的程序或内核 \\
基址/界限 & 地址先加基址，并检查是否在连续范围内 & 能保护单段程序，但不适合稀疏堆栈和共享库 \\
分页虚拟地址 & 虚拟页号查页表，得到物理页号和权限 & 每页可独立映射、共享、撤销、懒分配和保护 \\
\bottomrule
\end{longtable}

\begin{lstlisting}
load from *p:
  virtual address comes from register
  MMU reads page table entry
  MMU checks present, R/W/X, user/supervisor
  memory system reads physical address
  CPU returns data or raises exception
\end{lstlisting}

页表项不是单纯的 \code{virtual -> physical} 映射。它还告诉硬件：这一页当前是否存在，用户态能不能访问，能不能写，能不能执行。硬件每次访问都按这些位检查；检查失败时，CPU 进入异常入口，把 fault address 和原因交给内核。

把一个 64 位虚拟地址拆开看，会更具体。页大小若是 4KiB，低 12 bit 是页内偏移，高位是虚拟页号。页表翻译只替换页号部分：虚拟页号查到物理页框号，低 12 bit 原样保留。于是 \code{0x7f00_1234} 这样的虚拟地址可以映射到任意物理页框的偏移 \code{0x234} 处。OS 分配和回收的基本单位通常是页，因为硬件权限和翻译也是按页生效。

\begin{lstlisting}
virtual address:
  VPN = high bits of address
  offset = low 12 bits for a 4KiB page

translation:
  PTE[VPN] -> PPN + permission bits
  physical address = PPN || offset
\end{lstlisting}

这也解释了为什么同一个物理页可以出现在多个进程的虚拟地址里。共享库代码页可以被多个进程映射为只读可执行；\code{fork} 后的匿名页可以先让父子进程共同映射为只读，等某一方写入时触发 COW；内核也可以把同一块设备 buffer 映射到不同地址，用不同权限约束访问。虚拟地址提供的是命名空间，物理页提供的是实际存储，页表把二者在某个时刻绑定起来。

一次 \code{load rax, [p]} 实际上有三条常见路径。第一条是快路径：TLB 里已经有 \code{p} 所在虚拟页的翻译和权限，权限允许，cache 命中，数据很快返回。第二条是中间路径：TLB 没有翻译，但页表项存在且权限允许，硬件 page walker 读取页表，填入 TLB，再继续访问 cache。第三条是异常路径：页表项不存在或权限不允许，硬件无法直接完成这条指令，必须保存 fault 证据并进入内核。

\begin{longtable}{p{0.18\linewidth}p{0.38\linewidth}p{0.32\linewidth}}
\toprule
路径 & 硬件怎么走 & OS 是否接管 \\
\midrule
TLB 命中 & 虚拟页直接得到物理页号和权限，继续访问 cache & 不接管，程序感觉只是一次普通读取 \\
TLB miss 但 PTE 有效 & page walker 读页表，检查权限，填 TLB & 通常不接管，只是变慢 \\
PTE 不存在 & page walker 发现无有效映射 & 接管，内核查 VMA 后决定是否补页 \\
权限失败 & PTE 存在但写/执行/用户权限不允许 & 接管，可能 COW、发信号或杀进程 \\
cache miss & 翻译和权限都通过，但数据不在近端 cache & 不接管，继续向下级 cache 或 DRAM 请求 \\
\bottomrule
\end{longtable}

这里最容易混的是 TLB miss 和 page fault。TLB miss 只是硬件缓存没有命中；page fault 是硬件查到“这次访问不能按当前页表直接完成”。前者通常不进入内核，后者一定要给内核一个可解释现场。OS 修改页表、做 COW、撤销映射、切换地址空间时，真正要维护的就是这条证据链。

\begin{pdfdiagram}
  \node[os state, minimum width=23mm] (reg) at (0,12mm) {寄存器\\虚拟地址};
  \node[os compute, minimum width=25mm] (agu) at (34mm,12mm) {地址生成\\base+index+off};
  \node[os memory, minimum width=26mm] (tlb) at (70mm,12mm) {TLB\\最近翻译};
  \node[os memory, minimum width=29mm] (ptw) at (70mm,-12mm) {page walker\\读取页表};
  \node[os control, minimum width=27mm] (perm) at (106mm,12mm) {权限检查\\R/W/X U/S};
  \node[os memory, minimum width=29mm] (cache) at (143mm,12mm) {cache\\物理 line};
  \node[os memory, minimum width=29mm] (dram) at (143mm,-12mm) {内存控制器\\DRAM};
  \node[os control, minimum width=25mm] (fault) at (106mm,-12mm) {异常证据\\addr+reason};
  \draw[os strong bus] (reg) -- (agu);
  \draw[os strong bus] (agu) -- node[os label, above]{虚拟页号} (tlb);
  \draw[os bus] (tlb.south) -- node[os label, right]{miss} (ptw.north);
  \draw[os strong bus] (ptw.east) -| node[os label, below right]{PTE} (perm.south);
  \draw[os strong bus] (tlb.east) -- node[os label, above]{PPN+权限} (perm.west);
  \draw[os strong bus] (perm) -- node[os label, above]{物理地址} (cache);
  \draw[os bus] (cache.south) -- node[os label, right]{miss} (dram.north);
  \draw[os feedback] (perm.south) -- node[os label, right]{失败} (fault.north);
\end{pdfdiagram}
\diagramnote{一次 load/store 不是“CPU 直接读内存”，而是地址生成、TLB/page walk、权限检查、cache 和内存控制器串起来的结构路径；page fault 来自权限或映射检查，不是普通 cache miss。}

\begin{longtable}{p{0.2\linewidth}p{0.36\linewidth}p{0.32\linewidth}}
\toprule
事件 & 硬件含义 & OS 是否立刻参与 \\
\midrule
PTE not present & 页表项当前无有效物理页 & 进入内核缺页处理 \\
权限不允许 & 写只读页、用户访问内核页、执行 NX 页 & 进入内核，可能 COW、发信号或杀进程 \\
机器检查 & 硬件发现不可忽略的物理错误 & 内核记录、隔离、降级或 panic \\
\bottomrule
\end{longtable}

由此可见，OS 内存管理远不止“分配内存”四个字。页表权限给进程隔离提供硬件根；缺页异常让懒分配、文件映射和 COW 成立；权限异常让内核能拦住非法访问。先有这条隔离路径，之后才能讨论“怎样让它更快”。

页表检查解决了安全问题，但每次访问都完整读页表会太慢。于是硬件增加翻译缓存，常叫 TLB，translation lookaside buffer。TLB 保存最近的虚拟页到物理页翻译结果。TLB miss 只是“翻译缓存里没有”，硬件可以自己走页表；page fault 是“页表或权限告诉硬件这次访问不能直接完成”，CPU 必须进入内核。两者不能混同。

把缺页处理接到前面的提交边界上，就能理解懒分配为什么成立。用户程序第一次读 \code{p[0]}，硬件发现 PTE 不 present，保存 faulting RIP 和 fault address，进入内核。内核查到 \code{p} 落在合法 VMA 内，于是分配物理页、填 PTE、刷新必要的旧翻译，然后返回原指令。CPU 重新执行同一条 load，这次翻译成功，数据返回。用户程序看到的只是第一次访问慢一些；它没有看到中间的内核路径。

\begin{lstlisting}
lazy page allocation:
  user load p[0]
  hardware finds PTE not present
  kernel checks VMA and allocates a page
  kernel installs PTE and returns
  CPU retries the original load
\end{lstlisting}

如果 fault 之前 \code{rax} 已经被写回，重试就会破坏程序语义；如果 faulting RIP 保存错了，内核也不知道该回到哪条指令。虚拟内存能把“没有物理页”伪装成“第一次访问时再准备”，依赖的正是硬件把访存失败停在提交前。

数据访问本身也需要加速。主内存比执行单元慢得多，若每次 load/store 都等完整内存延迟，流水线会经常停住。于是硬件在核心附近放更小、更快的副本，叫 cache。cache 以 cache line 为单位保存最近访问过的物理内存内容。

cache 不是为了改变内存语义，而是为了利用局部性。若程序刚读过地址 \code{0x1008}，它很可能很快再读附近地址 \code{0x1010}；硬件一次从下级内存搬来一整条 cache line，例如 64 字节，后续相邻访问就能在 L1 cache 中完成。若程序反复遍历数组，cache 能显著减少访问 DRAM 的次数；若程序随机跳跃访问很大的内存区域，cache 命中率下降，执行单元就会更多等待。

\begin{longtable}{p{0.18\linewidth}p{0.38\linewidth}p{0.32\linewidth}}
\toprule
访问现象 & 硬件解释 & OS 或程序后果 \\
\midrule
顺序扫描数组快 & 相邻元素落在同一条或相邻 cache line & 预读、批量处理和局部性友好 \\
链表随机跳慢 & 每个节点可能在不同 cache line，甚至不同页 & 指针追逐造成 cache/TLB miss \\
任务迁移后变慢 & 新 CPU 的 L1/L2 和 TLB 没有旧任务热状态 & 调度器要考虑亲和性 \\
共享锁很慢 & 锁所在 cache line 在多个 CPU 间转移所有权 & per-CPU、分片和队列锁减少抖动 \\
\bottomrule
\end{longtable}

TLB 和 cache 常被一起提到，但它们缓存的东西不同。TLB 缓存“虚拟页到物理页以及权限”的翻译结果；cache 缓存“某个物理 cache line 的数据内容”。一次 load 通常先需要翻译，再按物理地址访问数据 cache。TLB miss 不等于数据 cache miss；数据 cache miss 也不等于 page fault。把这三者区分开，阅读内核日志和性能计数时才不会混乱。

\begin{longtable}{p{0.2\linewidth}p{0.36\linewidth}p{0.32\linewidth}}
\toprule
内存子系统部件 & 它实际做什么 & OS 关心点 \\
\midrule
L1/L2/L3 cache & 缓存物理 cache line，减少访存延迟 & 伪共享、锁竞争、cache 热度、DMA 同步 \\
TLB/page walker & 缓存或读取页表翻译结果 & shootdown、PCID、大页、COW 权限变化 \\
内存控制器 & 把物理地址事务送到 DRAM channel/rank/bank & NUMA、带宽、ECC 错误、内存热插拔 \\
NUMA 互联 & 连接多个 socket 或内存节点 & 线程亲和性、first-touch、页迁移 \\
内存类型机制 & 标记可缓存、不可缓存、write-combining 等属性 & MMIO 不能按普通 RAM 缓存，framebuffer 需要特殊属性 \\
\bottomrule
\end{longtable}

\begin{longtable}{p{0.2\linewidth}p{0.36\linewidth}p{0.32\linewidth}}
\toprule
事件 & 硬件含义 & OS 是否立刻参与 \\
\midrule
L1 cache miss & 数据不在最近的 cache 中 & 通常不参与，只是变慢 \\
TLB miss & 翻译不在 TLB 中 & 通常不参与，page walker 读页表 \\
page fault & 页表或权限不允许这次访问直接完成 & 进入内核，修复或报错 \\
\bottomrule
\end{longtable}

DRAM 也不是一个平面数组。内存控制器要在 channel、rank、bank、row 之间调度请求；远端 NUMA 节点要经过互联；ECC 可能报告可纠正或不可纠正错误。OS 的页分配器、NUMA 策略和错误处理因此不是“性能优化附录”，而是现代内存系统的必要配套。

将一次 cache miss 继续展开，可以看到更多硬件状态：

\begin{lstlisting}
cache miss to DRAM:
  pick memory controller by physical address
  choose channel/rank/bank
  if target row is not open:
    precharge old row
    activate target row
  read cache line
  check or correct ECC
  return line to last-level cache
\end{lstlisting}

这对应两个常见现象。第一，连续访问相邻地址通常更快，因为它们可能落在相同或相近的 DRAM 行和 cache line 上。第二，NUMA 机器上“同样是内存”也有远近，线程在 CPU0 上运行却频繁访问 CPU1 所属节点的内存，会让延迟和互联流量上升。OS 的 first-touch、页迁移、任务亲和性，都是在处理这类硬件距离。

当一个执行流不够用时，系统会增加多个执行核心。单核机器上，同一时刻只有一个执行流真正运行；多核机器上，多个 CPU 同时读写共享内存，核心问题转为：谁先看到谁的写入？同一条 cache line 在多个核心里如何保持一致？两个 CPU 同时改一个计数器会不会丢更新？

先看最小失败例子。内存里的 \code{counter} 初值是 0，CPU0 和 CPU1 同时执行 \code{counter++}。这条语句至少包含三步：load 旧值、ALU 加一、store 新值。若两个 CPU 都先读到 0，再各自写回 1，最后结果就是 1，而不是 2。问题不在 C 语句写得短，而在硬件动作不是不可分割的一步。

\begin{longtable}{p{0.16\linewidth}p{0.32\linewidth}p{0.32\linewidth}p{0.12\linewidth}}
\toprule
时刻 & CPU0 & CPU1 & 内存结果 \\
\midrule
开始 & 准备读 \code{counter} & 准备读 \code{counter} & 0 \\
读旧值 & 读到 0 & 读到 0 & 0 \\
本地计算 & ALU 得到 1 & ALU 得到 1 & 0 \\
写回 & 写 1 & 写 1 & 1 \\
\bottomrule
\end{longtable}

多核硬件必须先解决两个层面的问题。第一层是“一致性”：同一个物理地址的 cache line 不能让多个核心随意写出互相矛盾的版本，写入所有权要在核心之间转移。第二层是“顺序”：一个 CPU 先写数据再写标志，另一个 CPU 看到标志时，是否一定能看到数据。这两个问题分别对应 cache 一致性协议、原子指令和内存屏障。OS 的锁、引用计数、等待队列和 RCU 都建立在这些硬件能力上。

\begin{pdfdiagram}
  \node[os compute, minimum width=24mm] (c0) at (0,14mm) {Core 0\\执行流 A};
  \node[os memory, minimum width=24mm] (l10) at (34mm,14mm) {L1 cache\\line X};
  \node[os compute, minimum width=24mm] (c1) at (0,-14mm) {Core 1\\执行流 B};
  \node[os memory, minimum width=24mm] (l11) at (34mm,-14mm) {L1 cache\\line X};
  \node[os memory, minimum width=32mm] (llc) at (82mm,0) {共享 cache/互联\\一致性目录};
  \node[os memory, minimum width=30mm] (mem) at (128mm,0) {内存控制器\\DRAM};
  \draw[os strong bus] (c0) -- node[os label, above]{load/store} (l10);
  \draw[os strong bus] (c1) -- node[os label, below]{load/store} (l11);
  \draw[os bus] (l10.east) -- node[os label, above]{共享/独占状态} (llc.west);
  \draw[os bus] (l11.east) -- node[os label, below]{失效/转移所有权} (llc.west);
  \draw[os strong bus] (llc) -- node[os label, above]{cache line} (mem);
\end{pdfdiagram}
\diagramnote{多核共享内存不是多个 CPU 同时改一个平面数组，而是每个核心先操作自己的 cache line 副本，再由一致性协议转移共享或独占所有权。}

\begin{lstlisting}
counter++ is not one hardware action:
  load counter
  add 1 in ALU
  store counter
\end{lstlisting}

为解决这种并发更新问题，多核硬件提供原子读改写和 cache 一致性协议，OS 在其上建立锁、引用计数、RCU、per-CPU 数据和调度队列。原子操作的关键不在语法，而在硬件确保某个 cache line 的所有权和更新不可被其他 CPU 插入。

\begin{lstlisting}
atomic_fetch_add(counter, 1):
  get exclusive ownership of counter cache line
  read old value
  write old value + 1
  keep other cores from interleaving another write
\end{lstlisting}

原子操作的成本来自硬件所有权转移。若很多 CPU 同时修改同一个变量，这条 cache line 会在核心之间来回迁移；每个 CPU 都要等待自己取得独占所有权。屏障则规定“发布对象字段”和“发布对象指针”的可见顺序。

把锁放进这个结构看，会更清楚。获取锁通常意味着对锁变量执行原子读改写，硬件要让当前 CPU 取得锁所在 cache line 的独占所有权。释放锁通常是一次带 release 语义的 store，要求临界区里对共享数据的写入先对其他 CPU 可见，再让锁看起来已经释放。另一个 CPU 获取锁时使用 acquire 语义，要求它看到锁已获取后，也能看到前一个持有者发布的数据。锁保护的不是源码括号，而是硬件可见顺序。

\begin{longtable}{p{0.18\linewidth}p{0.38\linewidth}p{0.32\linewidth}}
\toprule
同步动作 & 硬件必须保证什么 & 缺失后果 \\
\midrule
原子加 & 读旧值和写新值之间不能插入另一个写者 & 计数丢失、引用计数提前归零 \\
获取锁 & 获取者之后的读写不能跑到获取之前 & 进入临界区却看不到受保护数据 \\
释放锁 & 临界区写入不能跑到释放锁之后 & 别的 CPU 看到锁空闲却读到旧数据 \\
唤醒等待者 & 状态更新先可见，再通知远端 CPU 或线程 & 被唤醒者看不到唤醒原因 \\
\bottomrule
\end{longtable}

\begin{longtable}{p{0.22\linewidth}p{0.34\linewidth}p{0.32\linewidth}}
\toprule
多核硬件事实 & 直接问题 & OS 应对方式 \\
\midrule
每核私有 cache & 同一数据可能有多个副本 & 一致性协议、cacheline 对齐、per-CPU 数据 \\
原子读改写要独占 line & 热锁会让 line 在 CPU 间迁移 & queued spinlock、分片计数、减少全局锁 \\
store buffer 延迟可见 & 源码顺序不等于其他 CPU 观察顺序 & acquire/release、内存屏障、锁语义 \\
跨核通知靠 IPI & 远端 CPU 需要被打断配合 & TLB shootdown、reschedule IPI、call function \\
NUMA 访问不等价 & 远端内存更慢，远端设备 DMA 更绕 & 调度亲和性、页迁移、IRQ affinity \\
\bottomrule
\end{longtable}

因此，内核经常把“同一件事”拆成每 CPU 状态。全局计数器最直观，但每次更新都要抢同一条 cache line；每 CPU 计数器读起来复杂一些，却能让热写路径留在本地。硬件迫使 OS 在“读时汇总”和“写时争抢”之间做取舍。

只会计算和访问内存的机器仍然是封闭的。真正的计算机还要接入磁盘、网卡、键盘、显示、USB、GPU、传感器和安全芯片。设备通信的最小需求，是把一个字节交给外部设备。

设备不是函数，而是另一个按自己时钟前进的状态机。最小输出设备至少需要两个可观察位置：一个状态位告诉内核“现在能不能收字节”，一个数据槽接收要输出的字节。

可以先把设备想成一台小机器，而不是一个 C 对象。它有自己的状态寄存器、数据寄存器、控制逻辑和外部引脚。CPU 不能直接进入设备内部调用“发送函数”，只能通过总线读写设备暴露出来的寄存器窗口。CPU 写入数据寄存器，只是把一个命令放到设备边界上；设备什么时候真正把 bit 送到线缆、屏幕或控制器外部，由设备自己的状态机决定。

\begin{lstlisting}
minimal output device:
  status slot:
    READY = 1 means device can accept one byte
  data slot:
    writing a byte starts device-side output work
\end{lstlisting}

从结构上看，这个设备内部至少也有保存状态的寄存器。状态寄存器保存 \code{READY}、\code{BUSY}、\code{ERR}；数据寄存器保存 CPU 写进来的字节；设备自己的控制逻辑在另一套时钟下把数据送到外部世界；完成后再把 \code{READY} 置回 1，或在失败时置 \code{ERR}。CPU 看见的不是设备内部所有细节，而是 MMIO 暴露出来的这些状态窗口。

\begin{pdfdiagram}[scale=0.84]
  \node[os compute, minimum width=25mm] (cpu) at (0,8mm) {CPU\\load/store};
  \node[os control, minimum width=27mm] (mmio) at (38mm,8mm) {MMIO 地址窗口\\状态/数据槽};
  \node[os state, minimum width=27mm] (status) at (80mm,20mm) {状态寄存器\\READY/ERR};
  \node[os state, minimum width=27mm] (data) at (80mm,-3mm) {数据寄存器\\byte};
  \node[os compute, minimum width=28mm] (logic) at (118mm,8mm) {设备控制逻辑\\发送/完成/失败};
  \node[os control, minimum width=24mm] (event) at (154mm,8mm) {事件输出\\状态或 IRQ};
  \draw[os strong bus] (cpu) -- node[os label, above]{物理地址事务} (mmio);
  \draw[os bus] (mmio.east) |- node[os label, above left]{读} (status.west);
  \draw[os strong bus] (mmio.east) |- node[os label, below left]{写} (data.west);
  \draw[os strong bus] (data.east) -- (logic.west);
  \draw[os feedback] (logic.north) -- node[os label, above]{更新 ready/err} (status.east);
  \draw[os feedback] (logic.east) -- node[os label, above]{完成/失败} (event.west);
\end{pdfdiagram}
\diagramnote{设备寄存器不是普通变量，而是 CPU 和设备状态机之间的边界。CPU 写数据寄存器只是启动设备侧动作，完成和失败要靠设备后来更新状态或发出事件。}

\code{output\_one\_byte('A')} 因此可以具体化为：

\begin{lstlisting}
output_one_byte('A'):
  while load(DEVICE_STATUS).READY == 0:
    keep waiting or check timeout
  store(DEVICE_DATA, 'A')
\end{lstlisting}

这里的 \code{load} 和 \code{store} 形式上像访问内存，但目标不是普通 RAM，而是设备暴露出来的状态槽。CPU 的地址生成、权限检查和总线事务仍然存在；只是物理地址译码发现目标落在设备 BAR，而不是落在 DRAM。把设备状态槽放进地址空间，让 CPU 用普通 load/store 指令访问，这种方式叫 memory-mapped I/O，简称 MMIO。MMIO 的关键在于访问规则不同：读状态可能有副作用，写数据可能启动设备动作，访问顺序也不能随意缓存和重排。

这一步把前面的内存路径改了一点：虚拟地址仍先经页表翻译成物理地址，但物理地址不一定通往 DRAM。平台地址译码会判断这段物理地址属于普通内存、固件保留区，还是某个设备 BAR。若落在设备 BAR，请求会经过片上互联或 PCIe root complex 到设备，而不是进入内存控制器。于是同样是 \code{store} 指令，写普通内存的含义是“更新某个 cache line”，写 MMIO 的含义是“向设备寄存器发出一个有顺序的事务”。

\begin{longtable}{p{0.18\linewidth}p{0.38\linewidth}p{0.32\linewidth}}
\toprule
访问目标 & 结构路径 & 语义差异 \\
\midrule
普通 RAM & CPU store -> cache -> 内存控制器 -> DRAM & 可缓存、可合并、可稍后写回 \\
设备 BAR & CPU store -> 不可普通缓存映射 -> root complex/总线 -> 设备寄存器 & 可能启动动作，顺序和宽度由设备手册规定 \\
设备状态读 & CPU load -> 总线读请求 -> 设备返回状态寄存器 & 读到的是设备当前状态，可能有副作用 \\
\bottomrule
\end{longtable}

这段代码可以按硬件时序展开。假设 \code{DEVICE_STATUS} 和 \code{DEVICE_DATA} 都落在同一个设备 BAR 内，且这段地址被内核映射成不可普通缓存的 MMIO 区间。CPU 每次执行 \code{load(DEVICE_STATUS)}，都会形成一次到设备的读事务；设备返回当前状态寄存器的值。若 \code{READY=0}，CPU 继续循环或检查超时计数。若 \code{READY=1}，CPU 执行 \code{store(DEVICE_DATA, 'A')}，这次写事务到达设备数据寄存器，设备控制逻辑把它解释成“开始输出这个字节”。

\begin{longtable}{p{0.16\linewidth}p{0.42\linewidth}p{0.3\linewidth}}
\toprule
步骤 & 硬件路径 & 驱动必须记住什么 \\
\midrule
读状态地址 & CPU 发起 MMIO read，地址译码到设备 BAR & 读到的是设备当前状态，不是普通变量副本 \\
看到 busy & \code{READY=0}，设备状态机还没准备好 & 不能无限空转，要有超时或让出 CPU \\
看到 ready & \code{READY=1}，数据寄存器可接收 & 下一步写数据才有意义 \\
写数据地址 & CPU 发起 MMIO write，字节进入设备数据寄存器 & 写入顺序就是设备命令顺序 \\
设备开始工作 & 设备清 \code{READY}，按自己时钟输出字节 & CPU 不再控制每一个外部 bit \\
完成或失败 & 设备置 \code{READY} 或 \code{ERR}，也可能发 IRQ & 驱动必须重新读状态确认结果 \\
\bottomrule
\end{longtable}

这也解释了为什么驱动代码里常见 \code{readl}、\code{writel} 这类专门接口，而不是随便拿普通指针读写。普通内存读写可以被编译器和 CPU 缓存、合并或重排；MMIO 读写是和外部状态机对话。把 \code{DEVICE_STATUS} 缓存在普通变量里，循环就可能永远看不到设备后来置回的 \code{READY}；把两个 MMIO 写重排，设备可能先收到“启动”，后收到“描述符地址”，请求就会乱掉。

\begin{pdfdiagram}
  \node[os compute, minimum width=25mm] (cpu) at (0,12mm) {CPU\\load/store};
  \node[os control, minimum width=28mm] (decode) at (37mm,12mm) {地址译码\\RAM 或 MMIO};
  \node[os memory, minimum width=27mm] (ram) at (78mm,23mm) {普通内存\\可缓存};
  \node[os control, minimum width=27mm] (bar) at (78mm,0mm) {设备 BAR\\不可普通缓存};
  \node[os state, minimum width=26mm] (status) at (118mm,10mm) {状态寄存器\\READY/ERR};
  \node[os state, minimum width=26mm] (data) at (118mm,-12mm) {数据寄存器\\写入启动动作};
  \draw[os strong bus] (cpu) -- (decode);
  \draw[os bus] (decode) -- node[os label, above]{物理地址落在 RAM} (ram);
  \draw[os strong bus] (decode) -- node[os label, above]{物理地址落在 BAR} (bar);
  \draw[os bus] (bar.east) -- (status.west);
  \draw[os strong bus] (bar.east) |- (data.west);
\end{pdfdiagram}
\diagramnote{MMIO 是把设备寄存器放进地址空间，但它不是普通 RAM；同样的 load/store 到了设备 BAR 后，语义变成读取状态或向设备状态机下命令。}

\begin{longtable}{p{0.22\linewidth}p{0.34\linewidth}p{0.32\linewidth}}
\toprule
需求增长 & 推出的机制 & OS 立刻遇到的问题 \\
\midrule
读设备是否 ready & 设备状态寄存器 & 轮询不能无限等，必须有超时和错误证据 \\
交给设备一个字节 & 设备数据寄存器 & 写入顺序就是命令顺序，不能按普通内存优化 \\
设备太慢 & 中断或可睡眠等待 & 线程不能一直空转，要能睡眠和被唤醒 \\
设备有很多请求 & 队列和 completion & 完成事件必须能找回原请求 \\
数据太大 & 描述符和设备侧搬运 & buffer 归属和生命周期必须明确 \\
\bottomrule
\end{longtable}

如果设备很快，轮询状态位的成本尚可接受；如果设备很慢，CPU 一直转圈等待会浪费执行资源。随之产生的需求是“设备完成时能通知内核”。硬件可以让设备或中断控制器发出 interrupt request，简称 IRQ；CPU 在指令边界进入中断入口，内核再读取设备状态，确认到底哪个请求完成或失败。中断不是设备调用内核函数，它只是让 CPU 进入一个受控入口，真正的解释仍由驱动完成。

从轮询到中断，本质上是把“谁负责反复查看状态”换了一种方式。轮询时，CPU 主动一遍遍读状态寄存器；中断时，驱动先提交请求并让当前线程睡眠，设备完成后制造一个硬件事件，CPU 进入中断入口，驱动再读状态寄存器或 completion 队列，最后唤醒等待者。完成事实仍然来自设备状态或 completion，不来自“中断已经发生”这件事本身。

\begin{longtable}{p{0.18\linewidth}p{0.38\linewidth}p{0.32\linewidth}}
\toprule
等待方式 & 硬件动作 & 适用边界 \\
\midrule
忙轮询 & CPU 不断发 MMIO read，看状态位变化 & 延迟低，但占用 CPU \\
带超时轮询 & CPU 周期性读取状态，同时检查时间边界 & 防止坏设备让内核永远等 \\
中断等待 & 设备完成后投递 IRQ，CPU 进入中断入口 & 省 CPU，但需要处理中断顺序和丢失事件 \\
队列 completion & 设备写 completion 项，驱动按 tag 找回请求 & 适合并发请求和批量 I/O \\
\bottomrule
\end{longtable}

如果一次只输出一个字节，CPU 自己写数据槽尚可接受；如果要读写磁盘块、网络包或显存区域，逐字节搬运会大量占用执行资源。随之产生的需求是“让设备按一批描述符自己搬大块数据”。这种直接内存访问叫 DMA，direct memory access。DMA 的风险也随之出现：设备不再只是被 CPU 读写状态槽，它也能主动读写内存。因此内核必须建立映射、限制设备能访问哪些 buffer，并在完成后撤销映射。约束设备访问范围的硬件叫 IOMMU。

从单字节设备过渡到队列设备，是因为“状态位 + 数据槽”一次只能表达一个很小的动作。磁盘读一个 4KiB 块、网卡接收一个包、NVMe 同时处理上千个请求，都需要把“要访问哪里、长度多少、方向是什么、完成后找谁”写成结构化记录。这个记录就是 descriptor。多个 descriptor 排成 submission queue；设备处理完后把结果写到 completion queue。doorbell 寄存器的作用不是传数据，而是告诉设备队列尾指针已经推进，可以去普通内存里读 descriptor 了。

\begin{longtable}{p{0.15\linewidth}p{0.43\linewidth}p{0.3\linewidth}}
\toprule
单字节模型 & 队列模型中的对应物 & 为什么需要升级 \\
\midrule
ready 状态位 & submission/completion queue 的 head、tail 和 phase & 多个请求不能只靠一个 ready bit 表达 \\
数据槽 & 内存中的 descriptor 和 buffer & 大块数据不能逐字节写 MMIO \\
写入启动 & 写 doorbell 通知设备读队列 & 命令和数据分离，减少 MMIO 次数 \\
完成状态 & completion 项或 MSI-X 中断 & 完成必须带 tag，才能找回具体请求 \\
\bottomrule
\end{longtable}

\begin{lstlisting}
modern queued device path:
  CPU maps device registers as MMIO
  driver allocates queue and DMA buffers
  CPU writes descriptors in RAM
  CPU writes MMIO doorbell
  device DMA-reads descriptors
  device DMA-writes data or completion
  device sends MSI-X interrupt
  kernel handler finds the completed request
\end{lstlisting}

由此形成驱动模型。驱动要发现设备、保留资源、映射 MMIO、分配 IRQ、设置 DMA mask、加入 IOMMU domain、管理队列、处理 reset 和热插拔。若没有 IOMMU，设备 DMA 错地址可能覆盖任意内存；若 completion 前释放 buffer，迟到 DMA 会写坏新对象；若中断处理顺序错，可能丢中断或中断风暴。

把一条现代队列请求走完整，会看到更多必须保持的顺序。驱动先在普通内存里准备 descriptor，里面放 DMA 地址、长度、方向和 request tag；再确保 descriptor 和 buffer 内容已经对设备可见；然后写 MMIO doorbell。doorbell 不是数据本身，它只是告诉设备“队列尾部推进了，请来读”。设备随后 DMA 读取 descriptor，经 IOMMU 把 IOVA 翻译成物理地址，再读写 buffer。完成后，设备写 completion 项，通常还会发 MSI-X 中断。内核中断处理函数不能只看 vector 就认为请求完成，它还要读 completion，按 tag 找回原请求，再做 cache 同步、unmap 和唤醒。

\begin{longtable}{p{0.14\linewidth}p{0.43\linewidth}p{0.31\linewidth}}
\toprule
阶段 & 结构路径 & 乱序或遗漏的后果 \\
\midrule
准备缓冲区 & CPU 写普通内存，必要时 pin 页或分配 DMA buffer & buffer 无效会导致 DMA 越界或写坏对象 \\
建立映射 & DMA API/IOMMU 建立 \code{IOVA -> PA} 和方向权限 & 设备可能访问不到页，或访问不该访问的页 \\
写描述符 & CPU 在提交队列写地址、长度、tag、flags & 设备读到旧 descriptor 会执行旧请求 \\
敲门铃 & CPU 对 MMIO BAR 写队列尾指针或命令 & posted write 可能稍后才到设备 \\
设备 DMA & 设备经 IOMMU 读 descriptor、读写 buffer & 映射或方向错会触发 IOMMU fault 或数据损坏 \\
写完成项 & 设备把结果、tag、状态写回 completion ring & tag 复用过早会唤醒错请求 \\
发现完成 & CPU 进入 IRQ 或轮询到新 completion & 必须再查 completion，不能只信 vector \\
回收请求 & 驱动同步 cache、unmap、释放页并唤醒等待者 & 过早释放会遇到迟到 DMA \\
\bottomrule
\end{longtable}

这一整条链路里，普通内存、MMIO、IOMMU 页表、设备内部队列和中断控制器都参与了同一个请求。驱动 bug 往往不是“少写一行寄存器”，而是某个结构边界的所有权或顺序没有讲清楚：descriptor 还没对设备可见就敲 doorbell，completion 没确认就释放 buffer，unmap 后设备还拿着旧 IOTLB 翻译继续 DMA。

\begin{pdfdiagram}
  \node[os compute, minimum width=25mm] (cpu) at (0,18mm) {CPU/驱动\\填描述符};
  \node[os memory, minimum width=30mm] (sq) at (42mm,18mm) {内存队列\\descriptor};
  \node[os control, minimum width=27mm] (door) at (42mm,-7mm) {MMIO doorbell\\通知设备};
  \node[os compute, minimum width=27mm] (dev) at (86mm,6mm) {设备控制器\\队列状态机};
  \node[os memory, minimum width=28mm] (iommu) at (128mm,18mm) {IOMMU\\IOVA->PA};
  \node[os memory, minimum width=28mm] (buf) at (128mm,-7mm) {DMA buffer\\数据页};
  \node[os control, minimum width=26mm] (irq) at (0,-22mm) {MSI-X/IRQ\\完成事件};
  \draw[os strong bus] (cpu) -- node[os label, above]{普通内存写} (sq);
  \draw[os bus] (cpu.south east) -- node[os label, below left]{MMIO 写} (door.west);
  \draw[os strong bus] (door.east) -- node[os label, above]{启动} (dev.west);
  \draw[os strong bus] (dev.north east) -- node[os label, above]{DMA 读描述符} (iommu.west);
  \draw[os strong bus] (iommu.south) -- node[os label, right]{检查并翻译} (buf.north);
  \draw[os bus] (dev.south) -- ++(0,-18mm) -| node[os label, near end, above]{完成事件} (irq.east);
  \draw[os feedback] (irq.north) -- node[os label, left]{进入内核入口} (cpu.south);
\end{pdfdiagram}
\diagramnote{现代设备请求跨过普通内存、MMIO、设备队列、IOMMU、DMA buffer 和中断入口；驱动管理的是这条结构路径上的所有权和顺序，而不是单独写一个寄存器。}

\begin{longtable}{p{0.2\linewidth}p{0.36\linewidth}p{0.32\linewidth}}
\toprule
设备事务细节 & 它意味着什么 & 驱动必须怎么处理 \\
\midrule
MMIO read 有副作用 & 读状态可能清位或弹出队列 & 用专门访问函数，按手册读写 \\
MMIO posted write & 写 doorbell 返回不等于设备已处理 & 需要屏障或读回形成排序点 \\
MSI-X 是消息写入 & 设备可把中断投递到不同 vector/CPU & 配置队列亲和性和中断亲和性 \\
completion ring 复用槽位 & 同一个数组格会多轮使用 & phase、generation 或 ownership 位防止误读旧完成 \\
设备 reset 不取消物理事实 & 迟到 DMA 或旧中断仍可能出现 & quiesce、mask IRQ、等待 in-flight、撤销映射 \\
\bottomrule
\end{longtable}

PCIe 设备还涉及一个基本事实：CPU 对设备 BAR 的写入通常会变成 PCIe posted write，也就是写事务发出去后 CPU 可以继续执行；而 MMIO read 通常要等设备或 root complex 返回 completion。这个差异会直接影响驱动里的排序。

\begin{lstlisting}
MMIO write doorbell:
  CPU store to uncached BAR mapping
  root complex emits PCIe posted write
  device receives command later

MMIO read status:
  CPU load from BAR mapping
  root complex emits PCIe read request
  CPU waits for completion data
\end{lstlisting}

驱动并不只是“写几个寄存器”。它管理的是 CPU、设备、IOMMU、中断控制器和内存之间的并发协议。一次请求只有在 descriptor 可见、doorbell 有序、DMA 映射有效、completion 被确认、unmap 完成之后，才真正回到内核可释放状态。

当计算、内存、隔离、多核、设备和 DMA 都出现后，机器就不再是几个孤立部件，而是一整个平台。现代硬件通常包含多核 CPU、私有和共享 cache、TLB、内存控制器、NUMA 节点、PCIe root complex、IOMMU、中断控制器、NVMe、网卡、GPU、USB 控制器、TPM、固件和电源管理单元。它们每个都有寄存器、队列、权限和错误状态。

\begin{pdfdiagram}[scale=0.84]
  \node[os compute, minimum width=28mm] (core0) at (0,18mm) {CPU core 0\\寄存器/TLB/L1};
  \node[os compute, minimum width=28mm] (core1) at (0,-8mm) {CPU core 1\\寄存器/TLB/L1};
  \node[os memory, minimum width=30mm] (llc) at (42mm,5mm) {共享 cache\\一致性互联};
  \node[os memory, minimum width=31mm] (mc) at (84mm,18mm) {内存控制器\\DRAM/ECC};
  \node[os control, minimum width=31mm] (rc) at (84mm,-8mm) {PCIe root complex\\地址路由};
  \node[os memory, minimum width=29mm] (iommu) at (128mm,5mm) {IOMMU\\设备页表};
  \node[os compute, minimum width=26mm] (dev1) at (156mm,20mm) {NVMe/网卡\\DMA 队列};
  \node[os control, minimum width=26mm] (intc) at (156mm,-7mm) {中断控制器\\vector/EOI};
  \node[os control, minimum width=27mm] (fw) at (128mm,-28mm) {固件/ACPI\\发现与电源};
  \draw[os strong bus] (core0.east) -- (llc.west);
  \draw[os strong bus] (core1.east) -- (llc.west);
  \draw[os strong bus] (llc.east) -- (mc.west);
  \draw[os bus] (llc.east) -- (rc.west);
  \draw[os strong bus] (rc.east) -- (iommu.west);
  \draw[os strong bus] (iommu.east) -- (dev1.west);
  \draw[os feedback] (dev1.south) -- node[os label, right]{MSI-X} (intc.north);
  \draw[os feedback] (intc.west) -- node[os label, below]{投递到 CPU} (core1.south east);
  \draw[os bus] (fw.north) -- node[os label, right]{描述资源} (rc.south);
\end{pdfdiagram}
\diagramnote{现代平台是一张互联结构图。OS 管理 task、地址空间、驱动和恢复路径时，实际是在维护 CPU 状态、内存翻译、设备 DMA、中断投递和固件描述之间的契约。}

\begin{longtable}{p{0.2\linewidth}p{0.36\linewidth}p{0.32\linewidth}}
\toprule
现代硬件块 & 主要接口 & OS 要维护的契约 \\
\midrule
CPU core & ISA、异常、中断、控制寄存器 & task、trap frame、系统调用、调度 \\
cache/TLB/MMU & load/store、PTE、TLB entry & 地址空间、COW、shootdown、cache 属性 \\
互联/NUMA & cache line、内存节点、远端访问 & 亲和性、页分配、负载均衡 \\
PCIe 设备 & BAR、MSI-X、DMA queue & probe、driver、IRQ、request lifecycle \\
IOMMU & IOVA、domain、IOTLB & DMA 隔离、VFIO、unmap 安全 \\
存储设备 & submission/completion queue、flush/FUA & page cache、块层、fsync、崩溃恢复 \\
安全硬件 & NX、SMEP/SMAP、TPM、Secure Boot & 隔离、启动信任链、最小权限 \\
\bottomrule
\end{longtable}

这条链路至此形成闭环：计算需求对应 ALU；保存和选择下一步对应寄存器、程序计数器和 CPU；多程序共存对应特权、异常、MMU 和页表；访问加速对应 TLB、cache 和内存层次；多个执行流并行对应一致性、原子指令和跨核通知；设备通信对应 MMIO、中断、队列和 DMA；设备隔离和启动可信性则对应 IOMMU 和安全硬件。OS 要管理的，正是这些状态、权限、等待和失败边界。

真实硬件会失败：DRAM 有 ECC 错误，PCIe 有 AER 错误，NVMe 请求会超时，网卡可能 reset，CPU 可能收到机器检查异常，电源管理会让设备暂时不可访问。OS 必须覆盖这些非正常路径，才能在长期运行中保持可解释状态。

\begin{longtable}{p{0.2\linewidth}p{0.36\linewidth}p{0.32\linewidth}}
\toprule
硬件错误来源 & 硬件留下的证据 & OS 恢复方向 \\
\midrule
CPU machine check & 错误 bank、地址、严重级别 & 隔离页、杀进程、panic 或降级 \\
DRAM ECC & 可纠正/不可纠正错误地址 & page offline、内存巡检、故障统计 \\
PCIe AER & 设备、链路、事务错误码 & reset 设备、重建队列、上报驱动 \\
NVMe timeout & request tag、队列状态、控制器状态 & abort、reset controller、失败请求 \\
中断丢失或风暴 & pending/in-service/mask 状态异常 & mask、轮询、清设备状态、EOI 顺序修正 \\
电源状态切换 & suspend/resume、reset、clock 状态 & 保存/恢复寄存器，停止和恢复 DMA \\
\bottomrule
\end{longtable}

恢复路径的核心是不变量：停止新请求，找出 in-flight 请求，阻止设备继续 DMA，撤销或重建映射，清理 completion，恢复寄存器和队列，再决定重试还是向上层报错。文件系统、网络和块层保存 request tag、引用计数和超时状态，目的就是在硬件偏离正常路径时把系统拉回可解释状态。

把这些硬件块放回一个小程序里，主线会更清楚。考虑下面这段代码：

\begin{lstlisting}
long sum_and_output(int *a, long n, int fd) {
  long s = 0;

  for (long i = 0; i < n; ++i) {
    s += a[i];
  }

  unsigned char out = (unsigned char)s;
  write(fd, &out, 1);
  return s;
}
\end{lstlisting}

源码只有一个循环和一次输出，但硬件不会看到“遍历数组”和“写文件描述符”这两个高层动作。硬件先看到的是函数入口处的一组寄存器参数：\code{a} 是一个虚拟地址，\code{n} 是循环上界，\code{fd} 是一个整数。循环变量 \code{i} 和累加器 \code{s} 可能放在通用寄存器里；每轮循环用地址生成单元算出 \code{a + i * 4}；load 指令把这个虚拟地址交给 MMU；ALU 把读回的整数加到 \code{s}；分支单元比较 \code{i} 和 \code{n}，决定下一 PC 回到循环开头还是继续向下执行。

\begin{longtable}{p{0.16\linewidth}p{0.43\linewidth}p{0.29\linewidth}}
\toprule
程序片段 & 硬件结构路径 & 关键边界 \\
\midrule
函数入口 & 调用约定把 \code{a}、\code{n}、\code{fd} 放进寄存器或栈，PC 指向函数第一条指令 & OS 切换任务时必须保存这些软件可见状态 \\
循环地址计算 & 寄存器文件给出 \code{a} 和 \code{i}，ALU/地址生成单元计算元素地址 & 地址仍是虚拟地址，还没有证明可读 \\
读取 \code{a[i]} & TLB 查翻译，必要时 page walk；权限通过后访问 cache、下级 cache 或 DRAM & TLB miss 只是慢；PTE 不存在或权限失败才进入内核 \\
累加 \code{s} & ALU 计算新和，寄存器写回在提交边界发生 & 异常或中断不能把半提交结果暴露给程序 \\
循环分支 & 比较单元设置条件，下一 PC 选择循环开头或退出路径 & 分支预测可提前取指，但猜错必须丢弃错误路径 \\
形成 \code{out} & 低 8 bit 被写到栈上或寄存器中的临时位置 & \code{\&out} 是用户虚拟地址，不是设备地址 \\
调用 \code{write} & 系统调用入口切到内核，内核检查 \code{fd} 和用户地址 & 用户参数不可信，必须转成内核对象和已验证地址 \\
设备输出 & 可能复制到内核 buffer，再走 MMIO 状态/数据寄存器，或走队列、DMA 和 completion & 返回值、设备完成和持久化边界不是同一个概念 \\
\bottomrule
\end{longtable}

现在把同一个循环放到几种硬件路径下对比。若数组刚被顺序访问过，\code{a[i]} 所在页的翻译可能在 TLB 中，目标 cache line 也可能在 L1 data cache 中。此时每轮循环主要消耗取指、译码、地址生成、L1 load、整数加法和分支预测。OS 不接管，用户程序只觉得循环很快。

若数组很大并且第一次扫描，路径会变长。第一次碰到某个虚拟页时，TLB 可能 miss；page walker 按当前页表基址读取多级页表，填入 TLB；目标数据不在 cache 时，内存请求继续走 L2、L3、内存控制器，可能还要跨 NUMA 互联到远端内存节点。硬件仍然不需要内核参与，因为页表允许访问；只是执行单元会等待更久，性能计数里会看到 TLB miss、cache miss 或远端内存访问。

若 \code{a[i]} 落在 \code{mmap} 后尚未装入的文件页上，事情才进入 OS。page walker 发现 PTE 不 present，CPU 保存 fault address、错误码和 faulting RIP，进入缺页异常入口。内核查 VMA，确认这个地址属于合法文件映射，于是分配物理页、向文件系统或块层读入内容、安装 PTE，并返回原来的 load 指令重新执行。用户代码仍然只是写了一次 \code{s += a[i]}，但第一次访问某页时，硬件和 OS 共同把“还没有物理页”补成了“现在可以读取”。

\begin{longtable}{p{0.18\linewidth}p{0.41\linewidth}p{0.29\linewidth}}
\toprule
同一轮循环的情况 & 实际路径 & 读者应该抓住什么 \\
\midrule
TLB 和 L1 都命中 & 地址翻译和数据都在核心附近，load 很快返回 & cache/TLB 是性能路径，OS 通常不出现 \\
TLB miss、PTE 有效 & 硬件 page walk 读取页表，填 TLB 后继续执行 & TLB miss 不等于 page fault \\
cache miss、权限通过 & 数据请求走下级 cache、DRAM 或远端 NUMA 节点 & 这是慢访问，不是权限失败 \\
PTE 不 present & CPU 进入缺页异常，内核根据 VMA 补页后重试 & 懒分配和文件映射依赖精确异常 \\
PTE 不允许写或执行 & CPU 在提交前进入保护异常 & COW、安全隔离和 NX 都靠权限位拦截 \\
定时器中断插入 & CPU 在指令边界保存现场进入内核，调度器可能换任务 & 循环没有主动让出 CPU，硬件定时器提供抢占边界 \\
\bottomrule
\end{longtable}

这个例子还说明“性能”和“语义”要分开读。分支预测、cache 命中、预取和乱序执行影响循环快慢；页表权限、异常现场、提交边界和系统调用入口决定程序是否合法、能否恢复、能否隔离。OS 可以利用性能结构，例如调度器尽量保持 cache 热度，页分配器考虑 NUMA；但 OS 的正确性不能建立在“这次应该会命中 cache”上。正确性必须建立在页表权限、特权级、精确异常、设备 completion 和可恢复现场这些硬边界上。

最后看 \code{write(fd, \&out, 1)}。这里的 \code{\&out} 不是设备数据寄存器地址，而是当前用户栈上的一个虚拟地址。系统调用入口只把这个地址作为不可信参数带入内核；内核必须先用文件表把 \code{fd} 转成 file 对象，再安全地从用户地址复制 1 个字节。如果这个 \code{fd} 指向普通文件，字节可能先进 page cache；如果它指向终端或串口，驱动可能轮询 MMIO 状态寄存器并写数据寄存器；如果它指向高速设备，内核可能把请求排到队列里，等待 DMA completion。一个统一的 \code{write} 接口背后，硬件路径取决于文件对象、内存属性、设备队列和完成语义。

从单字节输出推到现代平台之后，还要把几个关键路径压实。没有可恢复的执行现场，地址翻译、设备中断和 DMA 都无法接入 OS 主线；没有权限和事件证据，内核也无法判断一个动作是可重试、应报错，还是必须停止设备。

\section{二、CPU 执行、寄存器和特权入口}

上一节把固定数据通路推到了“能按指令选择动作”的机器。到了这一步，OS 最关心的不是 ALU 能算多快，而是两个更硬的事实：当前执行流能不能被暂停后恢复，用户代码能不能只从受控入口进入内核。进程切换、系统调用、缺页异常、中断和信号，全部压在这两个事实上。

先看一行普通 C 代码：

\begin{lstlisting}
x = a[i] + 1;
\end{lstlisting}

从源码看，它像一个整体动作；从硬件看，它早已被拆成若干条指令：取出 \code{a} 和 \code{i}，计算地址，读取 \code{a[i]}，加一，再把结果写回 \code{x}。如果读取 \code{a[i]} 时缺页，失败点不是“这一行 C 代码中间某处”，而是某条 load 指令。硬件必须告诉内核：哪条指令失败，访问哪个地址，失败原因是什么，哪些更早的结果已经提交，哪些更晚的猜测结果必须丢弃。

\begin{lstlisting}
possible instruction-level shape:
  addr = a + i * sizeof(a[0])
  tmp  = load [addr]       // may fault here
  tmp  = tmp + 1
  store [x], tmp           // may fault here too
\end{lstlisting}

把这件事再具体一点。假设有一个函数：

\begin{lstlisting}
int bump(int *a, long i, int *x) {
  int t = a[i] + 1;
  *x = t;
  return t;
}
\end{lstlisting}

按照常见 64 位调用约定，进入函数时可能是：\code{rdi} 保存 \code{a}，\code{rsi} 保存 \code{i}，\code{rdx} 保存 \code{x}。若 \code{a = 0x600000}，\code{i = 3}，\code{x = 0x700000}，并且 \code{int} 占 4 字节，那么 \code{a[i]} 的虚拟地址就是 \code{0x600000 + 3 * 4 = 0x60000c}。CPU 不知道这是数组第 3 个元素，它只看到一次地址生成、一次 load、一次整数加法和一次 store。

\begin{longtable}{p{0.16\linewidth}p{0.38\linewidth}p{0.34\linewidth}}
\toprule
机器动作 & 具体状态变化 & 失败边界 \\
\midrule
读参数寄存器 & \code{rdi=0x600000}，\code{rsi=3}，\code{rdx=0x700000} & 参数值本身来自调用者，硬件不懂 C 类型 \\
地址生成 & 地址生成单元计算 \code{0x60000c} & 这一步通常只是得到一个 bit 模式，还没检查页表 \\
读取数组元素 & MMU 检查 \code{0x60000c}，cache 或内存返回 4 字节 & 若页不存在或不可读，异常发生在写回结果之前 \\
整数加一 & ALU 把读出的值加 1，并产生候选结果 & ALU 结果还没有改变 \code{x} 指向的内存 \\
写回 \code{*x} & MMU 检查 \code{0x700000}，store 把结果写入目标地址 & 若目标页不可写，异常发生在 store 提交前 \\
返回值 & 返回寄存器保存同一个结果，下一 PC 回到调用者 & 调用者只看见已提交的寄存器和内存状态 \\
\bottomrule
\end{longtable}

这张表的重点是“同一行 C 代码”在机器里有多个可观察边界。若数组读取缺页，\code{*x} 还不能被修改；若写 \code{*x} 时触发 COW，数组读取和加法可以已经完成，但目标 store 还不能提交；若定时器中断正好打在加法和 store 之间，内核保存的现场必须足以让 CPU 将来回到 store 之前继续执行。读机器级程序时，要把源码动作不断压到寄存器、地址、cache、页表和提交边界上，而不是停在“执行这一行”。

这个例子先排除一个常见误解：CPU 不在执行 C 语句，而是在执行指令流。程序计数器指向下一条指令；通用寄存器保存整数、地址和临时值；状态寄存器保存条件码、中断使能、当前特权级等控制位；栈指针指向当前调用栈。OS 所谓“保存现场”，保存的就是这组软件可见状态。若少保存一个用户程序还会使用的寄存器，恢复后程序就会在看似随机的位置算错。

\begin{longtable}{p{0.2\linewidth}p{0.34\linewidth}p{0.34\linewidth}}
\toprule
硬件状态 & 在执行中的作用 & OS 关心点 \\
\midrule
程序计数器 & 指向下一条或当前异常指令 & 中断、异常、调度后从哪里继续 \\
栈指针 & 指向当前栈顶 & 用户栈和内核栈必须分开管理 \\
通用寄存器 & 保存参数、地址和临时值 & 系统调用参数、返回值、上下文保存 \\
状态寄存器 & 保存特权级、中断位、标志位 & 返回用户态前必须恢复正确权限 \\
控制寄存器 & 指向页表、控制地址翻译和中断 & 只有内核能修改关键控制状态 \\
浮点/向量寄存器 & 保存浮点和 SIMD 计算状态 & 状态很大，切换、信号和惰性保存要谨慎处理 \\
\bottomrule
\end{longtable}

从结构上看，CPU 同时维护两类状态。第一类是软件可见的架构状态：通用寄存器、RIP、RFLAGS、栈指针、页表基址和特权级。第二类是软件不可直接命名的内部状态：流水线寄存器、重排序缓冲、store buffer、分支预测器、cache 和 TLB。OS 切换任务时保存第一类状态；第二类状态通常由硬件自然清空、保留或按控制指令失效。这个区分解释了为什么“上下文切换”不是把 CPU 内部所有东西都倒出来保存。

取指、译码和执行不是记忆用的分类，而是硬件实现上的自然分工。取指单元只关心 RIP、指令 cache 和下一条地址；译码单元把二进制指令解释成内部控制信号；执行单元使用 ALU、分支单元、load/store 单元或浮点/SIMD 单元完成动作。把这些职责分开，CPU 才能流水化、预测和并行发射。

\begin{lstlisting}
front end:
  fetch bytes at RIP
  predict next RIP
  decode bytes into operations

back end:
  read registers
  execute ALU/branch/load/store
  access cache or memory
  commit architectural state
\end{lstlisting}

OS 关心这种划分有三个原因。第一，异常通常在执行或访存阶段才发现，但必须回报给架构层一个精确 RIP。第二，指令 cache、数据 cache 和 TLB 的冷热会影响系统调用、上下文切换和缺页性能。第三，分支预测和乱序执行虽然不改变正常架构语义，却会影响安全隔离和侧信道防护。

\begin{pdfdiagram}
  \node[os control, minimum width=30mm] (fe) at (0,12mm) {前端\\取指/预测/译码};
  \node[os state, minimum width=29mm] (sched) at (42mm,12mm) {调度/重命名\\微操作队列};
  \node[os compute, minimum width=31mm] (be) at (84mm,12mm) {后端执行\\ALU/branch/load/store};
  \node[os memory, minimum width=29mm] (mem) at (84mm,-12mm) {cache/TLB\\数据和翻译};
  \node[os state, minimum width=29mm] (rob) at (128mm,12mm) {提交队列\\按序提交};
  \node[os control, minimum width=29mm] (trap) at (128mm,-12mm) {异常入口\\精确现场};
  \draw[os strong bus] (fe) -- node[os label, above]{微操作} (sched);
  \draw[os strong bus] (sched) -- node[os label, above]{可执行操作} (be);
  \draw[os strong bus] (be.south) -- node[os label, right]{load/store} (mem.north);
  \draw[os strong bus] (be.east) -- node[os label, above]{完成结果} (rob.west);
  \draw[os feedback] (mem.east) -- node[os label, below]{fault 信息} (trap.west);
  \draw[os feedback] (rob.south) -- node[os label, right]{未提交时可丢弃} (trap.north);
\end{pdfdiagram}
\diagramnote{乱序执行可以改变内部完成顺序，但对 OS 暴露的异常必须发生在按序提交边界；否则缺页、信号和调试都无法可靠恢复。}

用户程序不能直接修改页表、关中断或访问任意设备寄存器。CPU 用特权级区分普通代码和内核代码。用户态执行特权指令会触发异常；系统调用是用户态主动进入内核态的受控入口；中断是设备或定时器异步打断 CPU 的入口；异常是当前指令执行失败产生的入口。

\begin{lstlisting}
user mode:
  execute normal instructions
  syscall -> controlled entry to kernel
  page fault -> exception entry to kernel
  timer interrupt -> asynchronous entry to kernel

kernel mode:
  validate context
  inspect object and permission
  update kernel state
  return to user or schedule another task
\end{lstlisting}

特权级解释了为什么系统调用不是普通函数调用。普通函数调用仍然在同一信任边界里；系统调用跨越用户/内核边界，必须切换栈、保存 trap frame、检查用户指针、处理信号和抢占，然后安全返回。

先把两者放在一起对比。用户态普通函数调用只是在同一个地址空间、同一个特权级、同一套用户栈上移动控制流：

\begin{lstlisting}
ordinary call:
  put arguments in registers or on user stack
  push or record return address
  jump to callee code in the same privilege level
  callee returns to caller
\end{lstlisting}

如果这个 callee 有 bug，它仍然只是用户进程的一部分；它不能因为被叫作“库函数”就修改页表或访问设备 BAR。系统调用不同。系统调用指令是 ISA 规定的入口，CPU 只允许跳到内核预先配置的入口地址，并把返回用户态所需的 RIP、状态位和参数现场保存到内核能解释的位置。用户不能选择跳进内核任意一行代码，也不能带着用户栈让内核随便使用。

\begin{longtable}{p{0.2\linewidth}p{0.34\linewidth}p{0.34\linewidth}}
\toprule
动作 & 普通函数调用 & 系统调用 \\
\midrule
特权级 & 不变，仍在用户态 & 从用户态进入内核态 \\
入口地址 & 由程序中的 call 目标决定 & 由 CPU 控制寄存器或向量表限定 \\
栈 & 使用用户栈 & 入口要切到内核栈或受控入口栈 \\
参数可信度 & 同一进程内部约定 & 全部来自不可信用户寄存器和用户地址 \\
返回条件 & callee 按 ABI 返回 & 内核恢复受控现场，执行返回用户态指令 \\
失败含义 & 进程内部 bug 或异常 & 可能是权限拒绝、地址错误、对象错误或设备错误 \\
\bottomrule
\end{longtable}

以 \code{write(fd, buf, len)} 为例，用户态库函数最多只是把参数放到 ABI 规定的寄存器里，再执行系统调用指令。真正的边界发生在硬件入口处。

\begin{lstlisting}
user:
  rax = SYS_write
  rdi = fd
  rsi = buf
  rdx = len
  syscall

hardware:
  save return RIP and status
  switch to kernel privilege
  jump to syscall entry

kernel:
  build trap frame
  validate fd and user buffer
  copy or pin data according to path
  return value or -errno
\end{lstlisting}

常见误解是：“内核已经拿到了 \code{buf} 指针，因此可以直接使用。”实际上，\code{buf} 只是用户虚拟地址。内核必须按用户地址空间检查它，拷贝时还可能跨页，途中也可能缺页。系统调用边界的本质不是函数名，而是：不可信寄存器参数进入可信内核入口，内核要把它们逐个变成经过验证的内核对象。

\begin{pdfdiagram}
  \node[os state, minimum width=26mm] (ureg) at (0,18mm) {用户寄存器\\fd, buf, len};
  \node[os control, minimum width=28mm] (entry) at (42mm,18mm) {syscall 指令\\切特权级};
  \node[os state, minimum width=29mm] (kstack) at (84mm,18mm) {内核栈\\trap frame};
  \node[os control, minimum width=29mm] (check) at (126mm,18mm) {内核检查\\fd 和权限};
  \node[os memory, minimum width=29mm] (copy) at (126mm,-8mm) {用户地址访问\\copy/pin};
  \node[os memory, minimum width=29mm] (pt) at (84mm,-8mm) {当前页表\\用户映射};
  \draw[os strong bus] (ureg) -- node[os label, above]{不可信参数} (entry);
  \draw[os strong bus] (entry) -- node[os label, above]{保存返回点/状态} (kstack);
  \draw[os strong bus] (kstack) -- node[os label, above]{构造内核上下文} (check);
  \draw[os strong bus] (check.south) -- node[os label, right]{访问 buf} (copy.north);
  \draw[os bus] (copy.west) -- node[os label, below]{仍受页表约束} (pt.east);
\end{pdfdiagram}
\diagramnote{系统调用入口只把 CPU 带到内核特权级，并不会把用户指针变成可信内核对象；参数还必须通过页表和内核对象表验证。}

要让特权级真正有用，CPU 不能只提供一个“当前模式”标志。它至少要提供五类硬件支持：受保护的控制寄存器、受控的陷入入口、返回用户态的指令、异常原因记录、以及对内存访问的权限检查。

\begin{longtable}{p{0.22\linewidth}p{0.34\linewidth}p{0.32\linewidth}}
\toprule
硬件能力 & 解决的问题 & 缺失后果 \\
\midrule
模式位/特权级 & 区分用户代码和内核代码 & 用户可执行所有敏感指令 \\
陷入向量表 & 限制用户只能从固定入口进内核 & 用户可跳到内核任意位置 \\
异常原因寄存器 & 告诉内核为何进入 & 无法区分缺页、非法指令和系统调用 \\
内核栈切换 & 避免使用不可信用户栈处理内核逻辑 & 用户栈损坏可破坏内核入口 \\
页表权限检查 & 限制用户访问地址范围和权限 & 进程可读写内核和其他进程 \\
\bottomrule
\end{longtable}

“把库函数写得安全一点”不能代替内核边界。库函数仍然运行在用户态，用户可以绕过库直接发指令。只有 CPU 在每条敏感路径上强制检查，OS 才能把不可信程序约束在用户态。

异常、陷入和中断来源不同，但入口结构相似：CPU 停下当前控制流，保存返回所需状态，切换到受信任入口，把原因交给内核，内核处理后决定返回、杀死任务、切换任务或 panic。

\begin{lstlisting}
hardware_entry(reason):
  save return RIP and status
  if crossing privilege boundary:
    switch to kernel stack
  set current privilege to kernel
  jump to vector[reason]

kernel_entry:
  save general registers
  build trap frame
  dispatch by reason
  decide return or schedule
\end{lstlisting}

入口路径必须极小心，因为它在系统最脆弱的位置运行。页表可能刚刚出错，用户指针不可信，中断状态敏感，栈也可能需要切换。许多架构把入口的一部分固定在汇编里，就是为了精确控制寄存器保存和返回协议。

CPU 如何停下当前程序并进入内核，已经有了硬件基础。进入内核后，几乎所有请求都会碰到地址：系统调用参数是地址，缺页异常是地址，设备寄存器也是地址。执行现场必须接上地址翻译和缓存，OS 才能继续判断“这个地址能不能碰、碰到的是 RAM 还是设备”。

\section{三、内存、设备和平台事务}

第二节解释了执行现场和特权入口，现在把一条具体地址继续往外走。用户程序里写下的地址，不会直接变成内存芯片上的位置；内核驱动里写下的设备地址，也不一定通向普通 RAM。同样是 load/store，硬件会先解释地址，再决定它落到普通内存、设备窗口、不可访问区域，还是一个必须交给内核处理的异常。

先从用户态的一次读取开始：

\begin{lstlisting}
char *p = mmap(file, length, PROT_READ);
char c = p[0];
\end{lstlisting}

\code{p} 是用户虚拟地址。CPU 发出的地址通常先交给 MMU。MMU 根据页表把虚拟页翻译成物理页，并检查 R/W/X、用户/内核、present 等权限。页表基址保存在控制寄存器中；TLB 缓存最近翻译，减少页表遍历成本。操作系统通过切换页表和设置权限，让每个进程看到自己的地址空间。

地址翻译的结构要分成三层看。第一层是 CPU 指令实际发出的虚拟地址，它来自寄存器、立即数和地址生成单元。第二层是硬件可读取的页表结构，它把虚拟页号连接到物理页框并附带权限。第三层是内核维护的 VMA 等软件结构，它说明某段虚拟区间是否合法、来自文件还是匿名内存、允许哪些访问。page fault 发生后，内核不是凭空分配页，而是在 VMA 允许的前提下修改页表，让硬件下一次访问能直接完成。

把一条 \code{mov rax, [rdi]} 放进这三层里看。寄存器 \code{rdi} 保存的是虚拟地址，例如 \code{0x7f00_0010}。地址生成单元把它作为本次 load 的虚拟地址交给 MMU。MMU 先用虚拟页号查 TLB；若命中，就拿到物理页号和权限；若未命中，就从当前页表基址开始做 page walk。page walk 读到的 PTE 如果 present 且允许用户读，硬件把物理页号和页内偏移拼成物理地址，再交给 cache。若 PTE 不存在或权限不允许，cache 根本不应该被当作成功路径继续访问，CPU 要记录 fault address 和错误码进入内核。

\begin{longtable}{p{0.17\linewidth}p{0.39\linewidth}p{0.32\linewidth}}
\toprule
阶段 & 硬件正在确认什么 & 失败时留下什么 \\
\midrule
地址生成 & \code{base + index * scale + offset} 得到虚拟地址 & 算术溢出通常不是地址 fault，本质仍是得到一个 bit 模式 \\
TLB 查询 & 最近是否已有该虚拟页翻译和权限 & TLB miss 只是慢，不等于异常 \\
page walk & 页表项是否存在，权限是否允许 & page fault address、访问类型、用户/内核、present 位状态 \\
cache/内存访问 & 物理 cache line 是否在近端，DRAM 是否返回 & cache miss 只是慢，机器检查才是硬错误 \\
\bottomrule
\end{longtable}

\begin{longtable}{p{0.22\linewidth}p{0.36\linewidth}p{0.3\linewidth}}
\toprule
概念 & 含义 & 常见错误理解 \\
\midrule
虚拟地址 & 程序使用的地址 & 误以为一定已经占用物理内存 \\
物理地址 & 内存控制器看到的地址 & 误以为用户程序能直接使用 \\
页表项 & 虚拟页到物理页和权限的记录 & 只当作普通 map，忽略权限 \\
TLB & 地址翻译缓存 & 修改页表后忘记失效旧翻译 \\
VMA & 合法虚拟区间的内核记录 & 误以为没有 PTE 就一定非法 \\
\bottomrule
\end{longtable}

页表是隔离的硬件根。没有页表权限，进程可以互相读写；没有用户/内核页区分，用户程序可以访问内核数据；没有 NX，数据页可能被当作代码执行。OS 的内存管理不是简单的“分配内存”，而是在性能、隔离、共享、懒加载和回收之间维护一组不变量。

为什么不让程序直接使用物理地址？早期简单机器可以这样做：地址 0x1000 对应某个内存芯片上的位置。但在多程序系统中，直接物理地址有三个根本缺陷。第一，程序彼此知道并可能覆盖对方地址。第二，程序必须按实际物理位置链接和装载。第三，内存碎片会让大程序难以找到连续空间。

虚拟地址把程序看到的地址和内存条上的物理地址分开。每个进程都可以以为自己从固定地址开始，有连续的栈、堆和映射区；内核用页表决定这些虚拟页实际落到哪些物理页，或者暂时不落到任何物理页。

\begin{longtable}{p{0.22\linewidth}p{0.34\linewidth}p{0.32\linewidth}}
\toprule
阶段 & 能力 & 局限 \\
\midrule
直接物理地址 & 实现最简单，裸机常见 & 无隔离，装载位置僵硬 \\
基址/界限 & 一个程序一段连续保护区 & 不适合稀疏地址和共享库 \\
分段 & 不同段不同基址和权限 & 外部碎片和模型复杂 \\
分页 & 固定大小页、独立权限和懒映射 & 页表/TLB 成本、内部碎片 \\
多级分页 & 稀疏地址空间节省页表内存 & page walk 更复杂，需 TLB 加速 \\
\bottomrule
\end{longtable}

分页之所以成为现代通用 OS 的核心，是因为它把“地址连续幻觉”和“物理页离散管理”解耦。进程堆可以在虚拟地址上连续增长，物理页却可以来自任意空闲页框；共享库可以映射到多个进程；\code{fork} 可以先共享物理页，等写入时再 COW。

页表项如果只保存 \code{vpn -> ppn}，只能回答“去哪里”，回答不了“能不能去”。真实页表项还包含权限、状态和性能位，因为硬件必须在每次地址翻译时直接做决定。

\begin{longtable}{p{0.18\linewidth}p{0.32\linewidth}p{0.38\linewidth}}
\toprule
位或字段 & 含义 & OS 用法 \\
\midrule
present/valid & 翻译是否有效 & 缺页、懒分配、换出 \\
PPN & 指向物理页框 & 页分配和共享的实际对象 \\
R/W/X & 读、写、执行权限 & 代码保护、COW、NX 防护 \\
U/S & 用户态是否可访问 & 内核页和用户页隔离 \\
accessed & 最近是否访问 & 页面回收近似 LRU \\
dirty & 是否写过 & 文件页 writeback、匿名页换出 \\
global/PCID & TLB 保留或地址空间标识 & 降低切换刷新成本 \\
\bottomrule
\end{longtable}

这些位让 OS 能把策略推迟到事件发生时。例如 \code{mmap} 文件时可以只建 VMA，不立刻读文件；第一次访问触发缺页，内核再读取对应页并安装 PTE。COW 也靠权限位实现：父子进程先共享只读页，写入触发保护异常，内核复制物理页并恢复可写权限。

把这件事落到 \code{mmap} 上，就能看到 OS 为什么要把策略推迟到第一次访问。把内存管理理解成“申请一块内存”过于粗略。实际路径是：建立虚拟区间，等待第一次访问，再由硬件把问题交给内核。

\begin{lstlisting}
char *p = mmap(file, length, PROT_READ);
char c = p[0];
\end{lstlisting}

\code{mmap} 返回后，\code{p} 这段地址通常有 VMA，不一定立刻有 PTE，也不一定已经读入物理页。VMA 是内核记录的“这段虚拟地址是否合法、权限是什么、背后对应哪个文件或匿名内存”；PTE 是硬件 page walk 能读取的“某个虚拟页当前是否映射到某个物理页，以及权限是什么”；TLB 是 CPU 对 PTE 结果的缓存。

\begin{pdfdiagram}
  \node[os state, minimum width=29mm] (vma) at (0,10mm) {VMA\\合法区间和策略};
  \node[os memory, minimum width=29mm] (pte) at (44mm,10mm) {PTE\\硬件可读映射};
  \node[os memory, minimum width=29mm] (tlb) at (88mm,10mm) {TLB\\翻译缓存};
  \node[os compute, minimum width=29mm] (load) at (132mm,10mm) {load 指令\\重试后完成};
  \node[os control, minimum width=29mm] (fault) at (44mm,-15mm) {page fault\\addr+error};
  \node[os memory, minimum width=29mm] (page) at (88mm,-15mm) {物理页\\文件页/匿名页};
  \draw[os feedback] (load.south) -- node[os label, right]{PTE 不 present} (fault.east);
  \draw[os strong bus] (fault.west) -- node[os label, above]{查合法性} (vma.east);
  \draw[os strong bus] (vma) -- node[os label, above]{允许后安装} (pte);
  \draw[os strong bus] (page.north) -- node[os label, right]{PPN} (pte.south);
  \draw[os bus] (pte) -- node[os label, above]{缓存翻译} (tlb);
  \draw[os strong bus] (tlb) -- node[os label, above]{下次命中} (load);
\end{pdfdiagram}
\diagramnote{VMA、PTE、TLB 不是并列术语。VMA 是内核策略，PTE 是硬件证据，TLB 是硬件缓存；缺页处理就是把合法的软件策略转成硬件可执行的页表项。}

\begin{longtable}{p{0.2\linewidth}p{0.38\linewidth}p{0.3\linewidth}}
\toprule
步骤 & 发生了什么 & OS 的判断 \\
\midrule
\code{mmap} 返回 & 内核建立 VMA，可能不安装 PTE & 这段地址合法，但不等于已有物理页 \\
CPU 执行 load & MMU 查 TLB，未命中后 page walk & 硬件先尝试自己翻译地址 \\
PTE 不 present & CPU 触发 page fault，记录 faulting address 和错误码 & 内核进入缺页处理，不是用户函数主动调用 \\
内核查 VMA & 地址落在合法 VMA 内且权限允许读 & 可以分配页或读取文件页 \\
安装 PTE & 页表项指向物理页，权限设为可读用户页 & 后续访问可由硬件直接完成 \\
返回原指令 & CPU 重新执行那条 load & 用户代码像是普通读取，只是第一次慢 \\
\bottomrule
\end{longtable}

若地址不在任何 VMA 内，内核通常给进程发 \code{SIGSEGV}；若 VMA 允许读但不允许写，写入会变成保护异常；若 PTE 改了但旧 TLB 还在，CPU 可能继续使用旧翻译，因此多核内核还要做 TLB shootdown。VMA、PTE 和 TLB 不是三个并列名词，而是一条访问路径上的三层证据。

地址和权限检查解决的是“能不能访问”，没有解决“访问快不快”。CPU 访问寄存器最快，其次是 L1/L2/L3 cache，再往后是内存和存储设备。cache 以 cache line 为单位搬运数据。多核系统中，同一个物理地址可能被多个 CPU cache 缓存，硬件一致性协议负责让普通内存读写最终保持一致。但设备 DMA、MMIO 和某些弱内存模型仍要求内核显式使用屏障和 cache 维护操作。

\begin{longtable}{p{0.18\linewidth}p{0.28\linewidth}p{0.42\linewidth}}
\toprule
层次 & 特点 & OS 影响 \\
\midrule
寄存器 & 每个 CPU 私有，最快 & 上下文切换必须保存可见寄存器 \\
L1/L2/L3 cache & 以 cache line 缓存内存 & 锁竞争、伪共享、迁移成本来自这里 \\
主内存 & 所有 CPU 共享 & 页分配、NUMA、DMA buffer 来自这里 \\
存储设备 & 持久但慢 & page cache、writeback、fsync 用来隐藏和约束它 \\
\bottomrule
\end{longtable}

缓存解释了许多 OS 设计选择。每 CPU runqueue 减少共享 cache line 抖动；per-CPU slab cache 减少全局锁；任务迁移会损失 cache/TLB 热状态；DMA 写入内存后，CPU 可能需要 invalidate cache 才能看到设备写入的新数据。

用两个很小的循环可以把 cache 和 TLB 的差异看得更清楚：

\begin{lstlisting}
long sum_array(int *a, long n) {
  long s = 0;
  for (long i = 0; i < n; ++i) {
    s += a[i];
  }
  return s;
}

long sum_stride(int *a, long n) {
  long s = 0;
  for (long i = 0; i < n; i += 1024) {
    s += a[i];
  }
  return s;
}
\end{lstlisting}

第一段顺序扫描数组。假设 cache line 是 64 字节，一个 \code{int} 是 4 字节，那么一条 cache line 能装 16 个连续元素。第一次访问 \code{a[0]} 可能 miss，硬件从下级 cache 或 DRAM 搬来包含 \code{a[0]} 到 \code{a[15]} 的整条 line；接下来访问 \code{a[1]}、\code{a[2]} 时通常命中同一条 line。地址翻译也有类似局部性：4KiB 页能装 1024 个 \code{int}，一次 TLB 翻译可以覆盖这一页内的连续元素。于是顺序扫描常见路径是“少量 cache/TLB miss 后，大量命中”。

第二段每次跳过 1024 个 \code{int}，步长正好是 4096 字节。每次访问都可能落到新页，也通常落到新 cache line。即使每次只读 4 字节，硬件仍会搬来一整条 cache line；但剩下 60 字节很可能用不上。若数组很大，TLB 也会频繁替换，因为每次访问都换一个虚拟页。程序源码只是把 \code{i += 1} 改成 \code{i += 1024}，硬件路径却从“line 内连续命中”变成“cache line 和 TLB 都不断换工作集”。

\begin{longtable}{p{0.18\linewidth}p{0.38\linewidth}p{0.32\linewidth}}
\toprule
访问模式 & 硬件看到什么 & OS 或程序后果 \\
\midrule
连续数组扫描 & 同一页内连续虚拟地址，同一 cache line 内连续物理字节 & cache line 利用率高，预取器容易工作，TLB 压力低 \\
大步长访问 & 每次跳到新 cache line，甚至新页 & 搬来的 line 大部分浪费，TLB 可能频繁 miss \\
链表追逐 & 下一个地址要等当前节点 load 返回后才知道 & 硬件难预取，乱序执行也难绕过依赖 \\
二维数组按行访问 & 内存布局和访问顺序一致，地址连续增长 & cache/TLB 局部性好 \\
二维数组按列访问 & 每次跨过一整行，步长可能很大 & cache line 利用率低，可能触发更多 TLB miss \\
\bottomrule
\end{longtable}

链表再差一步，因为它不仅可能跳到新 cache line，还存在真实数据依赖：

\begin{lstlisting}
long sum_list(Node *p) {
  long s = 0;
  while (p != nullptr) {
    s += p->value;
    p = p->next;
  }
  return s;
}
\end{lstlisting}

数组扫描时，下一次访问地址可以由循环变量提前算出；链表遍历时，下一次地址藏在当前节点的 \code{next} 字段里。当前节点没有从 cache 或 DRAM 返回之前，CPU 甚至不知道下一个节点在哪里。乱序核心可以先执行一些无关指令，却无法提前访问未知地址。OS 中许多数据结构会避免在热路径上做长链表追逐，或者把对象按 slab、per-CPU cache、批量数组和队列组织起来，原因就在这里。

cache line 还会把“看起来没有共享”的变量绑到同一条硬件路径上。cache 不是按单个字节搬运，而是按 cache line 搬运，常见大小是 64 字节量级。若两个 CPU 反复写同一 cache line 上的不同变量，硬件仍必须在两个 CPU 之间转移这条 line 的独占所有权。这叫伪共享：程序逻辑上没有共享同一个变量，硬件上却共享了同一个 cache line。

\begin{lstlisting}
struct BadCounter {
  atomic<uint64_t> cpu0_count;
  atomic<uint64_t> cpu1_count; // likely same cache line
};

struct BetterCounter {
  alignas(64) atomic<uint64_t> cpu0_count;
  alignas(64) atomic<uint64_t> cpu1_count;
};
\end{lstlisting}

OS 大量使用 per-CPU 数据，就是为了避免所有 CPU 争同一条 cache line。调度器每 CPU runqueue、网络每队列统计、slab 每 CPU 缓存，本质上都是把热写路径分散到局部状态，必要时再汇总。锁慢也常不是临界区代码慢，而是锁变量所在 cache line 在 CPU 间高速迁移。

以一次自旋锁交接为例。CPU0 修改受保护数据后释放锁，CPU1 随后获取同一把锁。锁变量所在 cache line 必须从 CPU0 的 cache 转到 CPU1，受保护数据所在的热 cache line 也可能要迁移。若所有 CPU 都抢同一把全局锁，性能下降不是因为“加锁语句多写了几行”，而是硬件互联在搬运 cache line 所有权。OS 把 runqueue、统计计数和内存分配缓存拆成 per-CPU 形式，就是为了让热写路径留在本地 CPU。

当多个核心和设备都能观察内存时，只写出源码顺序不够。内存屏障约束的不只是“C 语言语句顺序”。它同时约束编译器优化、CPU store buffer、cache 一致性可见顺序以及设备观察顺序。普通用户程序通常用高级原子语义间接表达；内核和驱动经常必须直接选择 acquire、release、full barrier、DMA barrier 或 MMIO barrier。

\begin{longtable}{p{0.22\linewidth}p{0.34\linewidth}p{0.32\linewidth}}
\toprule
场景 & 需要的顺序 & 典型工具 \\
\midrule
发布对象指针 & 对象字段先可见，指针后可见 & release/acquire \\
释放锁 & 临界区写入先可见，锁状态后释放 & unlock release 语义 \\
提交 DMA 描述符 & 描述符内存先对设备可见，再 doorbell & dma\_wmb/MMIO write barrier \\
读取设备完成 & 先看到 completion，再读取设备写入数据 & dma\_rmb/cache invalidate \\
\bottomrule
\end{longtable}

屏障不是越多越好。多余屏障会拖慢热路径；缺失屏障会造成难复现的并发错误。正确做法是先写出“谁发布什么，谁观察什么，观察到标志后必须看到哪些数据”，再选择足够但不过度的顺序约束。

CPU 和内存这条线解释了指令如何访问字节。真实 OS 还要和外部世界交换数据。外设不在 C 函数调用栈里，它们通过地址窗口、寄存器、队列和中断进入机器。普通内存访问因此必须继续延伸到设备访问。

设备通过总线连接到 CPU 和内存。很多设备把控制寄存器映射到物理地址空间，CPU 通过 MMIO 读写它们。MMIO 不是普通内存：读取可能有副作用，写入顺序可能被设备解释为命令序列，编译器和 CPU 不能随意优化这些访问。

总线描述的是读写事务的参与者和地址路由方式。CPU 访问普通内存时，请求通常走向内存控制器；CPU 访问某个 PCIe BAR 对应的物理地址时，请求走向 root complex，再被路由到具体设备；设备发起 DMA 时，请求反过来从设备进入 root complex，经 IOMMU 检查后到内存。

\begin{lstlisting}
CPU load/store:
  virtual address -> physical address
  physical address decoded by platform
  RAM range  -> memory controller
  MMIO range -> PCIe root complex -> device BAR

device DMA:
  device issues DMA address
  IOMMU translates and checks permission
  memory controller reads or writes DRAM
\end{lstlisting}

同样写成 \code{*addr = value}，普通内存和 MMIO 的语义并不相同。普通内存写入可以被 cache 合并、延迟、重排；MMIO 写入是在给设备状态机下命令，访问函数、内存属性和屏障都必须按设备手册来。

\begin{lstlisting}
driver_init(dev):
  writel(dev.mmio + RESET, 1)
  wait_until(readl(dev.mmio + STATUS) & READY)
  writel(dev.mmio + IRQ_ENABLE, RX_DONE | TX_DONE)
\end{lstlisting}

驱动必须给轮询加超时，因为设备可能坏掉。驱动也必须理解寄存器语义：有些状态位写 1 清除，有些寄存器读一次就弹出队列，有些寄存器要求先写高位再写低位。OS 不能把设备当作可靠函数调用，它面对的是异步硬件状态机。

最小 UART 驱动还能轮询一个状态位；高速设备通常不能这样工作。NVMe、网卡、virtio 这类设备更像一个队列状态机：CPU 在内存里填描述符，写 MMIO doorbell 通知设备；设备稍后 DMA 读取描述符、搬运数据、写 completion，再用中断或轮询告诉驱动。

\begin{lstlisting}
queue_device_submit(req):
  desc = reserve_submission_slot()
  desc.buffer = dma_address(req.buffer)
  desc.length = req.length
  desc.tag = req.tag
  dma_wmb()              // descriptor before doorbell
  writel(queue.tail, dev.doorbell)

queue_device_complete():
  cqe = read_completion_slot()
  req = find_request(cqe.tag)
  dma_unmap(req.buffer)
  wake_waiter(req)
\end{lstlisting}

在这段代码中，\code{desc} 位于普通内存，\code{doorbell} 位于 MMIO 区间，\code{buffer} 由设备通过 DMA 访问，\code{tag} 用于把异步 completion 找回原请求。若没有 \code{dma_wmb}，设备可能先看见 doorbell，再读到旧描述符；若 completion 前释放 buffer，设备迟到写入会破坏别的内核对象；若 tag 复用过早，完成事件可能唤醒错请求。

设备寄存器不是天然出现在内核变量里的。PCIe 设备刚被枚举时，OS 不能直接知道它的寄存器在哪里、需要几个中断、能做多少位 DMA。设备把一部分信息放在配置空间里。OS 读取 vendor id、device id、class code 判断它是什么设备；读取 BAR 判断设备希望暴露多大的 MMIO 区间；分配地址后再把这个区间映射进系统物理地址空间。

枚举过程可以理解成“把一块会响应总线事务的硬件，转成内核资源对象”。设备上电后并不会自动拥有一个内核虚拟地址；BAR 里描述的是它需要多大的地址窗口，系统固件或内核给这个窗口分配一段不和 RAM 冲突的物理地址范围。随后驱动用内核映射接口把这段物理 MMIO 区间映射到内核虚拟地址，访问属性设置成设备内存。到这一步，驱动手里的 \code{dev->mmio} 才是可以传给 \code{readl/writel} 的地址。

\begin{lstlisting}
PCIe BAR path:
  device exposes BAR size in config space
  OS allocates a physical MMIO range
  platform routes that range to the PCIe device
  driver maps the range as device memory
  readl/writel generate MMIO transactions to the BAR
\end{lstlisting}

\begin{longtable}{p{0.2\linewidth}p{0.36\linewidth}p{0.32\linewidth}}
\toprule
PCIe 信息 & 它说明什么 & OS 用它做什么 \\
\midrule
vendor/device id & 厂商和设备型号 & 匹配驱动或通用 class 驱动 \\
class code & 存储、网络、显示等设备类型 & 没有专用驱动时选择通用路径 \\
BAR & 设备寄存器或内存窗口大小 & 分配 MMIO 资源并建立内核映射 \\
capability list & MSI-X、PCIe、AER、电源管理等能力 & 配置中断、错误恢复和电源状态 \\
DMA mask & 设备能产生的地址位宽 & 决定 buffer 分配范围和 IOMMU 策略 \\
\bottomrule
\end{longtable}

这里的“资源”是硬边界，不是配置备注。两个设备不能占用同一段 MMIO；驱动不能访问未保留的 BAR；32 位 DMA 设备不能任意使用 4GiB 以上的物理页。设备枚举把硬件自描述信息转换成内核可管理的地址、中断和 DMA 约束。

在小型裸机程序里，外设地址通常写在芯片手册和头文件中；通用 OS 不能为每块主板硬编码所有设备。在本书的 x86-64 主线里，它需要从 UEFI/ACPI 和总线协议中发现硬件：ACPI 描述板级设备和电源关系，PCI/PCIe 配置空间描述可枚举设备，设备的 BAR 告诉内核 MMIO 区间，中断属性告诉内核如何接收事件，DMA mask 告诉内核设备能访问多少位地址。

\begin{lstlisting}
discover device:
  read firmware table or bus config space
  identify vendor/device/compatible string
  reserve MMIO region
  map MMIO into kernel virtual address
  allocate interrupt vector
  set DMA mask and IOMMU domain
  call matching driver probe()
\end{lstlisting}

这条路径解释了驱动 probe 为什么复杂。驱动不是拿到一个“设备对象”就能工作，它要把硬件描述转换成受内核管理的资源，并在任一步失败时回滚。MMIO、IRQ、DMA、clock、reset、power domain 和上层注册对象都属于生命周期的一部分。

probe 的顺序由资源依赖决定。内核必须先确认设备身份和资源边界，再映射寄存器，再建立中断和 DMA 约束，最后才允许设备开始工作。

\begin{longtable}{p{0.22\linewidth}p{0.36\linewidth}p{0.3\linewidth}}
\toprule
阶段 & 先做什么 & 过早或遗漏会怎样 \\
\midrule
发现设备 & 读取 ACPI/PCIe 信息，确认 vendor、class、BAR & 驱动可能绑定错设备或访问错地址 \\
保留资源 & 占用 MMIO 区间和 IRQ vector & 两个驱动可能同时操作同一硬件 \\
映射 MMIO & 用内核专门接口建立不可乱缓存的映射 & 普通指针访问会破坏设备顺序语义 \\
建立 DMA 边界 & 设置 DMA mask，加入 IOMMU domain & 设备可能访问不到 buffer，或越界访问内存 \\
启用中断/队列 & 初始化 ring、清状态、打开中断 & 设备可能在驱动未准备好时产生 completion \\
\bottomrule
\end{longtable}

设备完成发现、映射和初始化后，还需要把完成事件传回内核。设备不能直接调用内核函数返回，只能改变状态、写 completion，或者发出中断。事件从硬件回到 CPU 之后，定时器还给 OS 提供抢占和超时所需的时间边界。

中断让设备主动通知 CPU。定时器中断给 OS 提供抢占和超时能力；高精度时钟源提供时间戳；watchdog 检测系统是否长时间没有前进。没有定时器，任务死循环时系统无法抢回 CPU；没有中断，驱动只能轮询设备，浪费 CPU 并增加延迟。

\begin{longtable}{p{0.2\linewidth}p{0.34\linewidth}p{0.34\linewidth}}
\toprule
硬件机制 & OS 用途 & 失败表现 \\
\midrule
外设中断 & I/O completion、网络包、设备错误 & 中断风暴、丢中断、延迟尖峰 \\
定时器中断 & 时间片、超时、周期任务 & 任务霸占 CPU，超时不触发 \\
时钟源 & 时间戳、性能统计、调度账本 & 时间倒退或漂移导致判断错误 \\
watchdog & 检测软锁死和硬锁死 & 无响应时缺少故障证据 \\
\bottomrule
\end{longtable}

中断处理要短，因为它打断了当前执行流。复杂工作通常放到下半部、工作队列或内核线程。这个原则不是风格问题，而是延迟和可抢占性的硬边界。

旧式外设可能通过电平或边沿信号连到中断控制器。现代 PCIe 设备更常见的是 MSI 或 MSI-X：设备不是拉一根物理中断线，而是向一个特殊地址写一条消息。中断控制器接收消息，把它变成某个 vector，投递给某个 CPU。

注意这里的顺序：设备通常先把 completion 写入内存，再投递 MSI-X。MSI-X 本身也是设备发起的一次写消息，但目标不是普通数据 buffer，而是平台定义的中断投递地址。CPU 收到 vector 后，只知道“某个中断入口被触发”，还不知道哪个请求成功、哪个请求失败。驱动必须读取设备状态或 completion ring，根据 tag 找回请求对象。中断是门铃，不是结果本身。

\begin{lstlisting}
MSI-X interrupt delivery:
  device DMA-writes completion entry
  device writes MSI-X message
  interrupt controller receives vector
  target CPU takes interrupt at instruction boundary
  kernel handler reads device completion state
  kernel sends EOI when interrupt is handled
\end{lstlisting}

这条路径里有两个顺序不能混。第一，completion 数据必须先于中断可见，否则 CPU 进了中断处理函数却读不到完成项。第二，内核要先清设备侧原因，再对中断控制器 EOI；顺序反了可能立刻再次进入同一个中断，形成风暴。

\begin{longtable}{p{0.17\linewidth}p{0.39\linewidth}p{0.32\linewidth}}
\toprule
动作 & 结构含义 & 如果理解错了 \\
\midrule
设备写 completion & 结果被放到内存中的完成队列槽位 & 只看 IRQ 不看 completion 会丢失具体请求信息 \\
设备发 MSI-X & 设备向中断控制器投递消息 & MSI-X 不携带完整 I/O 结果 \\
CPU 进 vector & 硬件保存当前现场，跳到内核入口 & vector 可能共享或复用，不能当作唯一设备状态 \\
驱动读取完成队列 & 按 tag 匹配原请求并读取状态码 & tag 复用错误会唤醒错线程或释放错 buffer \\
EOI/重新打开中断 & 告诉中断控制器这次处理结束 & 过早 EOI 可能导致重入，过晚可能增加延迟 \\
\bottomrule
\end{longtable}

\begin{longtable}{p{0.2\linewidth}p{0.36\linewidth}p{0.32\linewidth}}
\toprule
中断概念 & 真实含义 & OS 风险 \\
\midrule
vector & CPU 入口编号，不等于设备对象 & 共享或复用时必须再查设备状态 \\
mask & 暂时禁止某路中断投递 & mask 后仍可能有 pending 事件 \\
pending & 事件已经到达但尚未处理 & 处理顺序错会丢事件或重复处理 \\
EOI & 告诉中断控制器本次处理完成 & 过早 EOI 可能造成重入或风暴 \\
affinity & 选择投递到哪个 CPU & 影响 cache 局部性、队列亲和性和负载 \\
\bottomrule
\end{longtable}

中断不是“设备调用内核函数”。设备最多只能制造一个硬件事件；内核还要从 vector 找到设备，从设备寄存器或 completion ring 找到具体请求，再把等待这个请求的线程唤醒。

定时器是抢占的硬件根。协作式系统依赖任务主动让出 CPU。只要当前任务不调用 \code{yield}，系统就无法切换。定时器打破了这个假设：硬件在未来某个时间点强制产生中断，CPU 无论当前运行哪个用户程序，都必须进入内核入口。内核因此获得重新选择任务的机会。

\begin{lstlisting}
timer interrupt:
  save current user/kernel context
  account elapsed time to current task
  run expired timers
  if current task exhausted slice:
    mark need_resched
  return path invokes scheduler if needed
\end{lstlisting}

定时器也支撑超时语义。睡眠、socket 超时、锁等待超时、TCP 重传、块设备请求超时、watchdog 都需要“未来某时刻再回来检查”。没有可靠时间源，OS 只能处理已经发生的事件，无法给等待设置边界。

用死循环看最清楚。用户程序可以写出永不主动让出 CPU 的代码：

\begin{lstlisting}
while (true) {
  // no syscall, no yield
}
\end{lstlisting}

若只有普通函数调用，内核永远拿不回控制权。定时器中断改变了这件事：内核在调度用户程序前设置下一次 clock event；时间到后，硬件强制进入中断入口；入口保存当前 RIP、状态寄存器和通用寄存器；内核给当前任务记账，发现时间片用尽，就把它放回可运行队列，选择另一个任务返回。即使用户程序没有主动配合，CPU 也能被抢回。

\begin{longtable}{p{0.2\linewidth}p{0.38\linewidth}p{0.3\linewidth}}
\toprule
时刻 & 硬件或内核动作 & 形成的 OS 概念 \\
\midrule
返回用户态前 & 内核设置下一次定时器事件 & 时间片有硬件边界 \\
用户死循环 & CPU 连续执行普通指令 & 任务当前占用 CPU \\
定时器到期 & CPU 异步进入中断向量 & 抢占入口出现 \\
入口保存现场 & trap frame 保存可恢复状态 & 任务可以被暂停 \\
调度器选择 & 当前任务未消失，只是回到 runqueue & 多任务共享 CPU \\
\bottomrule
\end{longtable}

“调度器能切任务”不是单靠算法实现的。算法只决定选谁，定时器和中断入口才提供“什么时候能强制重新选择”的硬件根。

OS 通常区分 clock source 和 clock event device。clock source 用来读取当前时间或单调计数；clock event device 用来安排下一次中断。前者回答“现在是多少”，后者回答“什么时候叫醒我”。

\begin{longtable}{p{0.22\linewidth}p{0.34\linewidth}p{0.32\linewidth}}
\toprule
类别 & 用途 & 错误影响 \\
\midrule
clock source & 时间戳、调度统计、超时比较 & 时间倒退、漂移、统计错误 \\
clock event & 周期 tick、hrtimer、抢占 & 超时不触发、任务不被抢占 \\
watchdog & 检测长期无进展 & 死锁或关中断过久无证据 \\
\bottomrule
\end{longtable}

tickless 内核不是没有定时器，而是不再固定频率无条件打断 CPU。内核计算最近的必要事件，把硬件定时器设置到那个时间点。这能降低空闲功耗和噪声，但要求调度、RCU、负载统计和 hrtimer 更精确地维护时间状态。

中断解决的是“设备或时间怎样通知 CPU”。但如果大块数据仍由 CPU 一字节一字节搬运，I/O 吞吐会受 CPU 搬运能力限制。PIO 由 CPU 搬运数据，简单但慢。DMA 让设备直接读写内存，CPU 只准备描述符、buffer 和地址映射。DMA 的难点是：设备访问的是 DMA 地址，不一定等同于 CPU 物理地址；buffer 在设备完成前不能释放；cache 可能需要 clean/invalidate；IOMMU 需要建立和拆除设备页表映射。

DMA 路径必须区分三种地址。用户进程拿到的是用户虚拟地址；内核管理页框时看到 CPU 物理地址；设备描述符里通常放 DMA 地址，也叫 IOVA 或总线地址。它们可能数值相同，也可能完全不同。驱动不能把用户指针直接塞给设备。

以 \code{write(fd, user_buf, 4096)} 走 DMA 为例，用户传进来的是用户虚拟地址。内核首先要把这段用户地址解析成一组物理页，并保证 I/O 期间这些页不能被释放或换出；然后通过 DMA API 为设备建立 IOVA 映射；最后把 IOVA 和长度写进 descriptor。设备看到的只有 descriptor 里的 DMA 地址，它不会理解用户进程的 \code{user_buf}，也不会自动遵守进程页表权限。这个转换链条是 DMA 安全性的核心。

\begin{longtable}{p{0.18\linewidth}p{0.38\linewidth}p{0.32\linewidth}}
\toprule
转换阶段 & 谁负责 & 不能跳过的原因 \\
\midrule
用户虚拟地址到页 & 内核按当前进程页表和 VMA 检查、pin 或复制 & 用户地址可能无效、跨页、权限不足或被并发撤销 \\
页到 DMA 映射 & DMA API/IOMMU 建立设备可见地址和方向权限 & 设备不能直接信任 CPU 物理地址，更不能拿用户指针 \\
DMA 地址到 descriptor & 驱动写入队列，并发布给设备 & 设备只读 descriptor，不读内核请求对象 \\
completion 到请求对象 & 驱动用 tag 反查原请求、同步 cache、unmap & 完成事件必须回收到正确 buffer 和等待者 \\
\bottomrule
\end{longtable}

\begin{longtable}{p{0.2\linewidth}p{0.36\linewidth}p{0.32\linewidth}}
\toprule
地址类型 & 谁使用它 & 为什么不能混用 \\
\midrule
用户虚拟地址 & 用户程序和系统调用参数 & 设备不理解进程页表，也不能信任用户权限 \\
内核虚拟地址 & 内核代码访问 buffer & 只对 CPU 有意义，设备看不到这个映射 \\
CPU 物理地址 & MMU 翻译后的内存地址 & 有 IOMMU 时未必等于设备可用地址 \\
DMA 地址/IOVA & 设备描述符中的地址 & 必须由 DMA API 建立和撤销映射 \\
\bottomrule
\end{longtable}

\begin{lstlisting}
submit_dma(req):
  req.page = pin_user_or_kernel_page(req.buffer)
  req.dma = dma_map(dev, req.page, req.len, TO_DEVICE)
  fill_descriptor(req.dma, req.len)
  memory_barrier()
  ring_doorbell(dev)
\end{lstlisting}

DMA 是高性能 I/O 的基础，也是内核缺陷的高风险区域。一个错误的 DMA 地址可能让设备覆盖内核对象；一个过早释放的 buffer 可能被设备迟到写入；一个缺失的 cache invalidate 可能让 CPU 读到旧数据。

DMA 的正确读法是“所有权转移”，不是“更快的 memcpy”。提交前 buffer 属于 CPU，CPU 可以写内容、填描述符、建 IOMMU 映射；提交后 buffer 属于设备，CPU 不能随意改、不能释放、不能复用；completion 到达并完成 unmap 后，buffer 才重新归 CPU。

\begin{longtable}{p{0.19\linewidth}p{0.39\linewidth}p{0.3\linewidth}}
\toprule
阶段 & 必须成立的不变量 & 破坏后的典型表现 \\
\midrule
map 前 & buffer 是有效页，方向和长度正确 & 设备读错数据或越界 \\
doorbell 前 & 描述符和 buffer 内容已经对设备可见 & 随机超时、设备读到旧命令 \\
DMA 期间 & CPU 不释放、不复用、不并发改写设备拥有的区域 & 迟到 DMA 写坏内存 \\
completion 后 & 先确认完成，再做 cache 同步和 unmap & CPU 读到旧数据或 IOVA 泄漏 \\
释放前 & 没有设备还持有该 buffer 的所有权 & use-after-free、IOMMU fault、随机崩溃 \\
\bottomrule
\end{longtable}

设备也需要访问边界。MMU 约束 CPU 访问内存，IOMMU 约束设备访问内存。没有 IOMMU 时，设备拿到的总线地址可能直接对应物理内存，驱动错误或恶意设备可以 DMA 到任意位置。有 IOMMU 后，内核为每个设备或设备组建立 DMA 地址空间，只映射允许访问的 buffer。

\begin{lstlisting}
dma_map(dev, page, len, direction):
  iova = allocate_device_virtual_address(dev.domain, len)
  iommu_map(dev.domain, iova, page.phys, len, permissions(direction))
  return iova

dma_unmap(dev, iova, len):
  iommu_unmap(dev.domain, iova, len)
  free_device_virtual_address(iova)
\end{lstlisting}

设备也有翻译缓存，通常称为 IOTLB 或类似结构。内核撤销 IOMMU 映射后，还要保证设备不会继续用旧翻译访问内存。这个要求和 CPU 修改页表后做 TLB shootdown 很像，只是参与者从 CPU 变成了设备。

\begin{lstlisting}
safe_dma_unmap(req):
  stop device from issuing new DMA for req
  observe completion or abort result
  iommu_unmap(req.iova)
  invalidate device/IOMMU translation cache if required
  release pages and request object
\end{lstlisting}

IOMMU 的代价是映射维护和 IOTLB 成本。高性能驱动会批量映射、复用 DMA buffer 或使用页池，减少每包 map/unmap。虚拟化场景中，IOMMU 还用于设备直通隔离：把设备 DMA 限制在某个虚拟机拥有的内存范围内。

用户页参与 DMA 时，还要区分 pin 和复制。当用户态提交异步 I/O 时，内核可以复制用户数据到内核 buffer，也可以 pin 用户页让设备直接 DMA。pin 的意思是这些物理页在 I/O 完成前不能被回收、迁移或换出。它避免复制，但把用户内存生命周期和设备生命周期绑在一起。

\begin{longtable}{p{0.2\linewidth}p{0.34\linewidth}p{0.34\linewidth}}
\toprule
策略 & 优点 & 代价 \\
\midrule
复制到内核 buffer & 生命周期简单，用户页可立即变化 & 多一次内存复制 \\
pin 用户页 DMA & 少复制，适合大 I/O & 页不能迁移/回收，completion 前必须保持有效 \\
共享 ring buffer & 系统调用少，批量提交完成 & 需要严格所有权和内存顺序协议 \\
\bottomrule
\end{longtable}

高性能 I/O 接口往往要显式注册 buffer、固定队列大小、区分 submit 和 completion。性能提升的代价，是把简单的同步复制转换成异步所有权协议。

用户页参与 DMA 时还要处理取消和进程退出。用户进程可以关闭 fd、收到信号、被终止，甚至把原来的虚拟地址 \code{munmap} 掉；但只要设备还可能写这个物理页，内核就不能把页还给普通分配器。可靠实现通常把请求对象、pinned page 列表、DMA 映射和 completion 绑定在一起，取消路径只是阻止新提交或标记请求，不等于立刻释放所有 buffer。

前面的讨论默认内核已经在运行，已经有页表、栈、日志和驱动框架。真实机器上这些条件不是天然存在的。最后一条硬件链路回到上电瞬间：第一条内核指令之前，谁初始化最低限度的 CPU、内存和设备状态。

机器上电后，CPU 不会直接运行 C 语言内核。固件先初始化基本硬件，bootloader 把内核镜像和启动参数放入内存，然后跳到内核入口。内核早期代码要设置栈、清 BSS、复制 data、建立早期页表、解析内存地图、初始化 console 和分配器。

\begin{lstlisting}
PowerOn
  -> Firmware
  -> Bootloader
  -> KernelEntry
  -> StackReady
  -> BssCleared
  -> EarlyConsole
  -> PageTablesReady
  -> SchedulerReady
  -> FirstUserProcess
\end{lstlisting}

启动阶段的约束最严格：还没有完整内存分配器，不能随意分配；还没有调度器，不能睡眠；还没有完整日志，错误证据有限；还没有完整设备模型，很多硬件不能访问。这些限制解释了内核初始化代码为何显得繁琐：早期代码不是在完整运行时环境中执行，而是在逐步制造这个环境。

早期内核不像普通 C 程序。普通 C 程序默认已经有栈、全局变量初始化、库函数、文件描述符、线程、堆和操作系统服务。内核入口处这些条件大多还不存在，内核必须自己逐步建立运行环境。

\begin{longtable}{p{0.22\linewidth}p{0.36\linewidth}p{0.3\linewidth}}
\toprule
早期阶段 & 能做什么 & 不能假设什么 \\
\midrule
刚进内核 & 少量汇编设置临时栈，确认启动参数 & 不能调用依赖完整栈和全局状态的复杂 C 代码 \\
清 BSS/定位镜像 & 建立全局变量的初始状态 & 不能假设静态对象已经自动可用 \\
早期页表 & 让内核虚拟地址能访问代码、数据和必要设备 & 不能随意访问未映射物理地址 \\
内存地图解析 & 找出可用 RAM、保留区、设备区和固件区 & 不能把所有物理地址都当可分配页 \\
早期 console & 尽早输出 panic 和调试证据 & 不能依赖完整驱动栈和文件系统日志 \\
分配器/调度器就绪 & 开始创建长期内核对象和任务 & 此前不能睡眠，也不能用普通锁等待别人唤醒 \\
\bottomrule
\end{longtable}

启动链把硬件概念收束到第一条内核指令：CPU 入口给出最初现场，页表决定哪些地址能用，内存地图决定哪些物理页能分配，MMIO console 提供早期证据，定时器和中断就绪后调度才有根，DMA 和完整设备模型要等内存、IOMMU 和驱动框架准备好之后才能安全启用。

\section{四、几个完整硬件路径大案例}

状态、ALU、CPU、页表、TLB、cache、MMIO、DMA、IOMMU、中断和多核一致性不能分开背。真正有用的读法，是把一次真实动作走穿：状态在哪里，数据走哪条通路，控制信号或权限怎样选择路径，失败或完成怎样留下证据。

\medskip
\noindent\textbf{第一条完整路径：同一条 \code{load rax, [p]} 为什么会有多条结果。}

假设用户态正在执行：

\begin{lstlisting}
rax = *p;
\end{lstlisting}

进入这条指令前，\code{p = 0x7f00_1234}。硬件先取指、译码，寄存器文件读出保存 \code{p} 的寄存器，地址生成单元得到虚拟地址 \code{0x7f00_1234}。从这一刻开始，路径会分叉。分叉点不是“是否走 cache”一个问题，而是依次经过地址合法性、翻译缓存、页表权限、数据缓存、内存控制器和错误报告。

\begin{longtable}{p{0.17\linewidth}p{0.43\linewidth}p{0.28\linewidth}}
\toprule
路径 & 硬件实际动作 & OS 看到什么 \\
\midrule
TLB 命中且 L1 命中 & TLB 给出物理页号和权限；物理地址命中 L1 data cache；数据返回执行单元；\code{rax} 在提交边界写回 & OS 不接管，用户程序只看到一次普通读取 \\
TLB miss 但 PTE 有效 & TLB 没有翻译；page walker 按页表基址读取多级页表；权限通过后填 TLB；再访问 cache & OS 通常不接管，只是这条 load 变慢 \\
cache miss & 翻译和权限都通过；L1/L2/LLC 没有目标 cache line；请求进入内存控制器或远端 NUMA 节点；返回整条 cache line & OS 通常不接管；性能计数会显示 cache/DRAM/NUMA 成本 \\
PTE 不存在 & page walker 发现页表项 not present；CPU 保存 fault address、error code 和 faulting RIP & OS 进入缺页处理，查 VMA 后补页、读文件页、分配匿名页或发信号 \\
权限失败 & PTE 存在但权限不允许，例如用户访问内核页、写只读页、执行 NX 页 & OS 接管，可能 COW、\code{SIGSEGV}、杀进程或 panic \\
机器检查 & 翻译和权限可能通过，但内存系统报告不可纠正 ECC 或总线错误 & OS 记录硬件错误，隔离页、杀进程、降级或 panic \\
\bottomrule
\end{longtable}

逐步看第一条快路径。TLB 命中说明硬件已经缓存了“虚拟页号到物理页号以及权限”的结果；L1 命中说明目标物理 cache line 已经在当前核心附近。此时 load 很快完成，\code{rax} 被写回，下一 PC 前进。OS 不应该介入，因为没有权限错误、没有缺页、没有设备事件。

第二条路径只是不在 TLB 中。硬件会自己走页表，这叫 page walk。page walk 读取的是内存里的页表项，所以它本身也可能访问 cache；若页表项在 cache 里，TLB miss 可能只是中等成本；若页表项也要去 DRAM，成本更高。但只要 PTE 有效且权限允许，这仍不是 page fault。把 TLB miss 当成缺页，是读 OS 时常见的概念错误。

第三条路径才是通常说的“不走 cache”或“cache miss”。准确说，它不是完全不经过 cache，而是先查近端 cache 未命中，然后向下级 cache、内存控制器、DRAM 或远端 NUMA 节点请求。返回的通常是一整条 cache line，而不是只返回 \code{rax} 需要的 8 字节。若 ECC 发现可纠正错误，硬件可能纠正后继续；若错误不可纠正，才会变成 OS 可见的硬件错误路径。

第四、第五条路径进入内核，是因为硬件无法把这条 load 当作当前页表下的合法动作完成。not-present 可能是懒分配、文件映射第一次访问、换出页或栈增长；权限失败可能是 COW 写保护、用户越界、执行不可执行页。内核能修复后重试，依赖前文讲过的提交边界：faulting 指令还没有把错误结果写入 \code{rax}，RIP 仍指向原 load。

\medskip
\noindent\textbf{第二条完整路径：同一条 \code{store [p], rax} 为什么不只是“写内存”。}

再看写入：

\begin{lstlisting}
*p = rax;
\end{lstlisting}

store 比 load 多一个容易忽略的事实：CPU 可以先把写入放进 store buffer，稍后再让它对其他 CPU、cache 层级或设备可见。对普通内存，这能提高吞吐；对锁和设备，这就必须用原子操作、release/acquire 或 MMIO 屏障表达顺序。

\begin{longtable}{p{0.2\linewidth}p{0.38\linewidth}p{0.3\linewidth}}
\toprule
写入场景 & 硬件路径 & OS 或并发后果 \\
\midrule
普通私有页写入 & 地址翻译通过；cache line 取得可写状态；store 先进入 store buffer，再更新本地 cache line，之后可能变成 dirty line & 用户程序正常前进；脏页以后可能被回写或换出 \\
写只读 COW 页 & PTE present 但不可写；CPU 在提交前触发保护异常 & 内核复制物理页，更新 PTE 后重试原 store \\
写共享 cache line & 当前 CPU 必须通过一致性协议取得该 line 的独占所有权；其他 CPU 的副本失效或降级 & 热锁、共享计数器和伪共享会让 line 在 CPU 间迁移 \\
原子读改写 & CPU 取得独占所有权，读旧值、算新值、写回期间不允许其他写者插入 & 引用计数、锁、队列指针才不会丢更新 \\
发布对象指针 & 数据字段写入和指针写入都可能经过 store buffer；release 屏障要求字段先可见，指针后可见 & 没有顺序约束时，其他 CPU 可能看到指针却看不到完整对象 \\
写 MMIO doorbell & 地址属性是设备内存；写事务进入总线或 PCIe posted write；返回不等于设备已经处理 & 驱动要用 \code{writel}、DMA 屏障或读回形成排序点 \\
\bottomrule
\end{longtable}

一个具体失败例子是发布队列节点：

\begin{lstlisting}
node->value = 42;
queue->tail = node;
\end{lstlisting}

源码顺序看起来是先写字段，再发布指针。但在硬件层面，两个 store 可能先进入 store buffer；另一个 CPU 若先看见 \code{queue->tail = node}，却还没看见 \code{node->value = 42}，就会读到未初始化内容。正确同步不是“相信 C 代码顺序”，而是用 release store 发布指针，用 acquire load 读取指针。release/acquire 约束的是其他 CPU 能观察到的硬件可见顺序。

写共享变量时还要考虑 cache line 所有权。假设两个 CPU 分别更新两个不同变量，但它们刚好落在同一条 64 字节 cache line。程序逻辑上没有共享变量，硬件却必须在两个 CPU 之间转移同一条 line 的独占所有权，这就是伪共享。OS 的 per-CPU 计数器、cacheline 对齐和分片锁，都是为了减少这类所有权迁移。

\medskip
\noindent\textbf{第三条完整路径：同样是地址访问，普通 RAM、MMIO 和 framebuffer 路径不同。}

CPU 指令形式上都可能是 \code{load/store}，但物理地址属性会改变语义。假设三个地址：

\begin{lstlisting}
normal_ptr      -> ordinary RAM
dev->status     -> MMIO status register
framebuffer_ptr -> write-combining device memory
\end{lstlisting}

这三个地址都要经过页表翻译，但 PTE 或平台内存类型会告诉硬件后续访问规则不同。

\begin{longtable}{p{0.2\linewidth}p{0.38\linewidth}p{0.3\linewidth}}
\toprule
目标 & 硬件访问规则 & 错用后的问题 \\
\midrule
普通 RAM & 可进入 cache，可预取，可合并写入，可延迟回写；一致性协议维护 CPU 间普通内存一致性 & 若拿它当设备寄存器，读写可能被缓存或重排，设备看不到真实命令顺序 \\
MMIO 状态寄存器 & 通常不可普通缓存；load 会变成总线读事务；读可能清状态位或弹出队列项 & 若缓存状态值，驱动可能永远看不到 ready 或错误位变化 \\
MMIO doorbell & store 变成设备命令；PCIe 上可能是 posted write，CPU 返回时设备未必已处理 & descriptor 未对设备可见就敲门铃，会让设备读到旧请求 \\
framebuffer/WC 区 & 写入可合并成更大的总线事务，提高顺序写带宽；读通常很慢或语义特殊 & 若按普通 RAM 随机读改写，性能和顺序都可能出问题 \\
DMA coherent buffer & CPU 和设备共享的内存区域，平台保证或由 API 维护可见性 & 忘记同步会导致 CPU 读旧数据或设备读旧描述符 \\
\bottomrule
\end{longtable}

这个案例说明，\code{cache} 只是问题的一部分。真正决定路径的是“地址翻译结果 + 内存类型 + 访问主体”。普通 RAM 的优化目标是让 CPU 快；MMIO 的目标是让设备按手册看到命令；framebuffer 的目标可能是顺序写吞吐；DMA buffer 的目标是让 CPU 和设备在明确的所有权边界上共享数据。OS 映射地址时必须把这些属性设对，否则同一条 \code{store} 指令会被硬件按错误规则执行。

\medskip
\noindent\textbf{第四条完整路径：一次 NVMe 读请求从系统调用走到 DMA 完成。}

再看一个大 I/O。用户执行：

\begin{lstlisting}
read(fd, user_buf, 4096);
\end{lstlisting}

\code{read} 不一定每次都碰设备。page cache 是内核用普通内存保存文件页内容的一层缓存。若目标文件页已经在 page cache 里，内核可以把这 4096 字节从内核页复制到用户 buffer；若不在，内核才需要向块设备提交 I/O，把数据读进 page cache 或目标页，再完成用户可见的读取。两条路径都叫 \code{read}，但硬件路径差别很大。

\begin{longtable}{p{0.18\linewidth}p{0.4\linewidth}p{0.3\linewidth}}
\toprule
路径 & 机器和内核动作 & 用户看到什么 \\
\midrule
page cache 命中 & 系统调用进入内核；内核找到已缓存文件页；CPU 通过普通 load/store 和 cache 把数据复制到 \code{user\_buf}；复制时仍要检查用户页权限，可能触发用户页缺页 & 没有设备 I/O，返回较快，但仍有系统调用和内存复制成本 \\
page cache 未命中 & 内核分配或选择页，提交块层请求；设备 DMA 把数据读入内存；completion 后等待者被唤醒；内核再把数据交给用户 buffer 或让映射页可用 & 当前线程可能睡眠，返回时间取决于设备、队列和调度 \\
用户 buffer 不可写 & 即使文件页命中，复制到用户地址时页表权限失败或地址非法 & 返回 \code{-EFAULT} 或产生相应错误，不应把设备路径误当根因 \\
设备错误 & cache 未命中后设备返回介质错误、超时或控制器错误 & 内核按块层和文件系统语义返回 I/O 错误 \\
\bottomrule
\end{longtable}

先看命中路径。CPU 从当前进程寄存器里拿到 \code{user\_buf}，系统调用入口把它带进内核；内核通过文件对象和页缓存索引找到目标文件页。接下来不是设备在搬数据，而是 CPU 在复制：读 page cache 页的 cache line，写用户地址对应的 cache line。写用户地址前，MMU 仍要按用户页表检查 \code{user\_buf} 是否可写；若用户 buffer 跨两个页，其中第二页不可写，复制可能在中途失败。内核必须保证返回值能表达“实际复制了多少”或返回错误，而不是假设一次 4096 字节复制必然整体成功。

未命中路径才会走到 NVMe。若数据不在 page cache，块层可能向 NVMe 设备提交读请求。这个路径同时经过系统调用入口、用户地址检查、页分配或 pin、DMA 映射、普通内存 descriptor、MMIO doorbell、IOMMU、设备 DMA、completion、MSI-X 中断和等待队列。

\begin{longtable}{p{0.15\linewidth}p{0.43\linewidth}p{0.3\linewidth}}
\toprule
阶段 & 硬件或内核动作 & 关键不变量 \\
\midrule
系统调用入口 & CPU 从用户态进入内核，保存返回 RIP、状态和寄存器参数 & \code{user_buf} 只是用户虚拟地址，不是设备可用地址 \\
检查用户 buffer & 内核按 VMA 和页表检查范围，必要时 fault-in 或 pin 页 & I/O 完成前相关页不能被释放、换出或迁移到不可访问状态 \\
建立 DMA 映射 & DMA API/IOMMU 建立 \code{IOVA -> PA}，附带读写方向权限 & 设备只能访问被授权的页和方向 \\
填写 submission entry & CPU 在普通内存队列写 opcode、LBA、长度、DMA 地址、request tag & descriptor 必须先对设备可见 \\
敲 doorbell & 驱动写 MMIO doorbell，通知设备队列 tail 前进 & doorbell 是设备命令，不是数据本身；posted write 需要排序意识 \\
设备取请求 & NVMe 控制器 DMA 读 submission entry，经 IOMMU 翻译 IOVA & IOMMU 或 IOTLB 错误会变成 DMA fault 或设备错误 \\
设备写数据 & 设备把磁盘数据 DMA 写入 buffer & CPU 不能提前释放或复用 buffer；非一致平台还要 cache invalidate \\
设备写 completion & 设备把 status、tag、phase 写入 completion queue，并投递 MSI-X & 中断只是通知，真正结果在 completion 项里 \\
内核完成请求 & IRQ handler 或轮询读取 completion，按 tag 找回请求，unmap，唤醒等待线程 & completion 前释放对象会导致迟到 DMA 或唤醒错请求 \\
\bottomrule
\end{longtable}

这条路径说明，“走不走 cache”无法单独解释 I/O。descriptor 在普通内存里，CPU 写它时会受 cache 和 store buffer 影响；doorbell 在 MMIO 区间，不能按普通内存缓存；设备 DMA 访问 buffer，要经过 IOMMU；completion 写回内存后，中断控制器把 MSI-X 投递给 CPU；CPU 再进入中断入口。任何一段顺序错了，表现可能是超时、数据损坏、中断风暴、IOMMU fault 或 use-after-free。

把这条读请求按硬件块拆开，能看出每个机制解决的不是同一个问题：

\begin{longtable}{p{0.18\linewidth}p{0.38\linewidth}p{0.32\linewidth}}
\toprule
硬件块 & 在这次 NVMe 读里负责什么 & 不能混淆的边界 \\
\midrule
TLB/MMU & 内核访问用户地址、队列和 MMIO 映射时，都要先把虚拟地址翻译成可访问的物理或设备地址范围 & TLB miss 只是翻译缓存未命中；用户 buffer 不合法才会走 fault 或错误返回 \\
CPU cache & CPU 写 submission entry、读 completion、复制用户数据时，会缓存普通内存里的 cache line & cache 命中只说明 CPU 访问普通内存快，不说明设备已经看到 descriptor \\
store buffer/屏障 & CPU 可能把普通内存写暂存在本地；屏障约束 descriptor 先于 doorbell 对设备可见 & 源码先写 descriptor 再写 doorbell，不自动等于设备观察顺序正确 \\
MMIO/BAR & doorbell 是设备寄存器，CPU 写它是在通知控制器读队列 & MMIO 不是普通 RAM，不能被普通缓存和随意合并替代 \\
DMA 引擎 & 设备按 descriptor 中的 DMA 地址读写 buffer 和 completion queue & DMA 不是 CPU memcpy，设备会在 CPU 之外主动访问内存 \\
IOMMU/IOTLB & 把设备看到的 IOVA 翻译成被授权的物理页，并限制方向和范围 & 它保护的是设备访问内存的边界，不是用户进程 load/store 的边界 \\
中断控制器/MSI-X & completion 写好后，把“有完成需要处理”的事件投递给某个 CPU & 中断只是提示，不携带完整请求结果；驱动仍要读 completion ring \\
request tag & 把异步 completion 对回原始内核请求对象和等待者 & tag 复用过早会把完成事件交给错误请求 \\
\bottomrule
\end{longtable}

因此，cache 只是路径中的一层。一次 I/O 还同时经过地址翻译、普通内存可见性、设备寄存器顺序、设备侧 DMA、IOMMU 权限、中断投递和请求匹配。优化或排错时只盯着 cache，容易把权限错误、顺序错误、设备超时和 completion 匹配错误都看成同一种慢访问。

\medskip
\noindent\textbf{第五条完整路径：一次 \code{write} 到 \code{fsync} 为什么才牵出持久化硬件。}

普通文件写入最容易被误解成“用户 buffer 里的字节马上写到磁盘”。实际情况通常不是这样。用户执行：

\begin{lstlisting}
write(fd, user_buf, 4096);
fsync(fd);
\end{lstlisting}

若 \code{fd} 是普通文件，\code{write} 的常见快路径是把用户数据复制到 page cache 中的文件页，并把这页标记为 dirty。dirty 的意思是“内存里的文件页比设备上的持久副本更新”。\code{write} 返回时，数据可能还只在 DRAM 里；后台 writeback 或之后的 \code{fsync} 才会把 dirty 页提交给块设备。于是 \code{write} 解决的是“内核已经接收这批字节”，\code{fsync} 才试图解决“这批字节和必要元数据已经进入设备承诺的持久化边界”。

\begin{longtable}{p{0.17\linewidth}p{0.43\linewidth}p{0.28\linewidth}}
\toprule
阶段 & 硬件和内核动作 & 关键边界 \\
\midrule
系统调用入口 & CPU 保存用户返回点和寄存器参数，进入内核 & \code{user\_buf} 仍是不可信用户虚拟地址 \\
复制用户数据 & 内核按用户页表读取 \code{user\_buf}，通过 CPU cache/load/store 复制到 page cache 页 & 用户页可能缺页或不可读，错误属于用户地址边界 \\
标记 dirty & 文件页在内存中变成新版本，并记录需要回写 & \code{write} 返回不等于设备已经持久化 \\
writeback 提交 & 内核把 dirty 页转成 bio/request，块层可能合并、排序或限流 & 文件页所有权从“只在内存中更新”转向“等待设备写入” \\
DMA 到设备 & 驱动建立 DMA 映射，填描述符，写 doorbell；设备读取内存中的数据并写入介质或控制器缓存 & buffer、descriptor、IOMMU 和 completion 顺序必须正确 \\
设备 completion & 设备报告写命令完成 & completion 可能只说明设备已接收或写入其缓存，是否断电不丢取决于设备缓存策略 \\
\code{fsync} & 内核等待相关 dirty 数据、日志或元数据提交，并在需要时发送 flush 或 FUA & 返回前必须满足文件系统和设备共同定义的持久化承诺 \\
\bottomrule
\end{longtable}

这里有两个不同层面的缓存。CPU cache 缓存的是普通内存里的 cache line，例如 page cache 页、描述符和 completion ring。存储设备也可能有自己的写缓存，可能在控制器 DRAM、SSD FTL 缓冲区或带断电保护的缓存中。CPU cache miss、page cache dirty、设备写缓存，这三个概念不能混在一起。CPU cache 是处理器近端副本；page cache 是内核用 DRAM 保存文件内容；设备写缓存是存储控制器为了吞吐暂存请求的内部状态。

\begin{longtable}{p{0.18\linewidth}p{0.4\linewidth}p{0.3\linewidth}}
\toprule
缓存层次 & 保存什么 & 断电或错误时的含义 \\
\midrule
CPU cache & 普通内存 cache line 的近端副本 & 硬件一致性通常保证最终写回 DRAM，但它不是文件持久化边界 \\
page cache & 文件页在内核 DRAM 中的副本 & \code{write} 后可能只有这里是新数据，断电会丢 \\
块层队列 & 等待提交、合并或调度的 I/O 请求 & 请求还没到设备，不能承诺持久化 \\
设备写缓存 & 控制器内部暂存的写入 & 若没有断电保护，completion 后仍可能在突然断电时丢失 \\
非易失介质或保护域 & NAND、磁盘介质，或有电容保护的控制器缓存 & \code{fsync} 想达到的通常是这个边界 \\
\bottomrule
\end{longtable}

持久化难点在于“写入顺序”也必须被保存。假设文件系统先写元数据说“文件长度已经增长”，但数据页还没真正落到持久介质，断电后就可能出现文件长度变大而内容是旧数据或空洞。文件系统因此会使用日志、写屏障、flush、FUA 或有序 writeback。FUA 可以理解成“这次写不要只停在易失写缓存里，完成前要到达设备承诺的持久边界”；flush 则是要求设备把之前已经接收的写缓存刷到持久边界。

\begin{lstlisting}
fsync_persistence_path:
  copy user bytes into page cache
  mark file page dirty
  submit data writeback
  wait for data write completion
  submit or commit filesystem metadata/journal
  issue flush or FUA if device write cache needs draining
  return success only after required completions
\end{lstlisting}

这个伪代码里的每一行都对应不同硬件证据。数据 write completion 说明那批数据请求有结果；journal completion 说明文件系统元数据的事务有结果；flush completion 说明设备承认之前的易失写缓存已经被推进到承诺边界。若设备谎报 flush，或者没有正确实现断电保护，OS 再严格也无法凭空制造持久性。OS 的职责是把请求顺序和设备命令发对，并在设备报告错误时向上层返回失败。

\begin{longtable}{p{0.2\linewidth}p{0.38\linewidth}p{0.3\linewidth}}
\toprule
断点 & 可能发生什么 & 恢复后看到的结果 \\
\midrule
\code{write} 返回后、writeback 前断电 & 新数据只在 DRAM page cache 里 & 文件可能仍是旧内容，除非另有同步保障 \\
设备写命令已完成、flush 前断电 & 数据可能在设备易失写缓存中 & 取决于设备是否有断电保护或是否要求 flush/FUA \\
数据持久化、元数据未持久化 & 数据块可能存在，但文件系统还没把它连接到文件名、长度或块映射 & 日志恢复可能丢弃这次变化，避免结构损坏 \\
元数据先持久化、数据未持久化 & 文件系统结构看似指向新数据，但内容可能旧或未定义 & 这是文件系统必须通过顺序和日志避免的危险路径 \\
\code{fsync} 成功返回后断电 & 设备和文件系统已经给出持久化完成证据 & 正常情况下应能恢复到包含这次同步写入的状态 \\
\bottomrule
\end{longtable}

这条路径把硬件链路再扩了一层：CPU 负责复制和提交请求，MMU 负责验证用户地址，page cache 负责吸收普通文件读写，块层负责排队和合并，NVMe 或其他控制器负责 DMA 和 completion，存储介质或设备保护域负责最终持久性。若把 \code{write} 理解成“直接写盘”，就无法解释为什么系统需要 dirty page、writeback、journal、flush、FUA、I/O scheduler 和错误上报。

\medskip
\noindent\textbf{第六条完整路径：改页表权限后，为什么还要让其他 CPU 配合。}

最后看一个纯内存管理的大案例。内核把某个虚拟页从可写改成只读，常见于 COW、\code{mprotect} 或撤销映射。

\begin{lstlisting}
make page read-only:
  update PTE writable bit from 1 to 0
  flush stale TLB entry on CPUs that may use it
  return after old writable translations cannot be used
\end{lstlisting}

修改内存里的 PTE 只是第一步。其他 CPU 的 TLB 里可能还缓存着旧的“可写”翻译；当前 CPU 的 store buffer 里也可能有更早的写入；另一个 CPU 可能正在用户态运行同一个进程。若不处理旧翻译，硬件仍可能按旧权限完成写入。

\begin{longtable}{p{0.18\linewidth}p{0.4\linewidth}p{0.3\linewidth}}
\toprule
阶段 & 硬件事实 & OS 为什么要处理 \\
\midrule
更新 PTE & 页表内存里的 writable 位被清掉 & 只改变了内存中的权威记录，没清掉 TLB 副本 \\
本地 TLB flush & 当前 CPU 丢弃该虚拟页旧翻译 & 防止当前 CPU 继续用旧权限写入 \\
远端 shootdown & 向可能运行该地址空间的 CPU 发 IPI，让它们失效旧 TLB entry & 防止其他 CPU 用旧 writable 翻译越过新权限 \\
等待确认 & 远端 CPU 完成失效后返回确认 & 内核才能认为旧权限已经不可再用 \\
后续写入 & 用户再次写该页时触发保护异常 & COW 或 \code{SIGSEGV} 才有可靠硬件证据 \\
\bottomrule
\end{longtable}

这个案例同时涉及页表、TLB、中断、跨核通知和调度。TLB shootdown 不是“刷新缓存”四个字那么简单：它要求 OS 知道哪些 CPU 可能持有这个地址空间，给它们发 IPI，远端 CPU 在安全点执行失效，再返回确认。若某个 CPU 正在关中断或长时间不可响应，权限切换的完成时间就会变长。COW、\code{mprotect}、进程迁移和安全隔离都压在这条硬件链路上。

\medskip
\noindent\textbf{第七条完整路径：一次 \code{write(fd, buf, 1)} 怎样从用户态走到设备。}

现在把开头的输出一个字节放进真实 OS 路径。用户态执行：

\begin{lstlisting}
write(fd, buf, 1);
\end{lstlisting}

若 \code{fd} 对应的是一个简单串口或早期 console，这个请求最后可能落到 MMIO 状态寄存器和数据寄存器上。注意这不是“用户调用驱动函数”。用户程序只能把参数放进寄存器，然后用系统调用指令进入内核；内核再把不可信参数转换成经过验证的内核对象和地址。

\begin{longtable}{p{0.15\linewidth}p{0.44\linewidth}p{0.29\linewidth}}
\toprule
阶段 & 硬件路径 & OS 必须确认什么 \\
\midrule
用户发起 & ABI 把 \code{fd}、\code{buf}、\code{len} 放进寄存器；系统调用指令触发受控入口 & 寄存器参数来自用户，不可信 \\
进入内核 & CPU 切换特权级，保存返回 RIP、状态寄存器和必要寄存器，改用内核栈 & trap frame 足够让内核将来返回用户态 \\
查 \code{fd} & 内核用当前 task 的文件表把小整数 \code{fd} 转成 file 对象 & \code{fd} 可能无效，或对象不支持写 \\
取用户字节 & \code{buf} 是用户虚拟地址；\code{copy\_from\_user} 会走用户页表、TLB 和 cache，也可能缺页 & 不能把用户指针直接当内核指针，更不能直接交给设备 \\
进入驱动 & 驱动拿到已经复制到内核侧的字节，获取设备锁，检查设备状态 & 多个写者不能同时打乱设备寄存器顺序 \\
读状态寄存器 & CPU 对 MMIO 状态地址发起 read，经地址译码、root complex 或片上总线到设备 & 读到的是设备当前状态，不是 cache 里的普通变量 \\
写数据寄存器 & CPU 对 MMIO 数据地址发起 write，设备状态机接收字节并清 busy/ready 状态 & 写入可能启动外部动作，不能按普通 RAM 重排 \\
返回用户态 & 内核设置返回值，系统调用返回路径恢复用户寄存器和特权级 & 返回成功只说明字节已提交给对应语义，不一定代表外设物理动作已经完全结束 \\
\bottomrule
\end{longtable}

这里同时有两条访存路径。读取 \code{buf[0]} 是普通用户内存访问：它要通过用户页表和 cache，可能命中 L1，也可能 TLB miss、缺页或权限失败。写设备数据寄存器是 MMIO 访问：它也从一条 store 指令开始，但地址属性把它送到设备寄存器，而不是送到普通 cache line。若把二者混成“读一个地址、写一个地址”，就会漏掉系统调用入口、用户地址验证、文件对象、驱动锁、MMIO 顺序和设备状态机。

失败路径也要按硬件边界拆开。系统调用入口失败，表现为非法入口或错误的寄存器状态；\code{fd} 失败，是内核对象表查不到；\code{buf} 失败，可能是页表权限或缺页处理不能完成；设备失败，可能是 MMIO 状态长期 busy、错误位被置起，或写后设备不再产生预期完成事件。每一类失败对应的证据不同，所以 OS 不能只返回一个模糊的“写失败”，它要尽量把错误定位到参数、地址、权限、设备状态或超时。

\medskip
\noindent\textbf{第八条完整路径：一次定时器中断怎样变成线程切换。}

再看没有用户主动调用的路径。当前 CPU 正在用户态执行普通代码，硬件定时器到期后投递中断。这个事件可能导致内核抢占当前线程，切换到另一个线程运行。

\begin{lstlisting}
timer expires:
  interrupt controller delivers vector
  CPU enters kernel interrupt path
  scheduler may choose another task
  return to selected task
\end{lstlisting}

这条路径重点不是 cache，而是“异步事件怎样在安全边界接入当前执行流”。定时器不是修改用户变量，也不是调用用户函数；它通过中断控制器让 CPU 在指令边界进入内核入口。内核随后把当前执行现场变成可保存对象，再决定是否继续运行原线程。

\begin{longtable}{p{0.15\linewidth}p{0.44\linewidth}p{0.29\linewidth}}
\toprule
阶段 & 硬件路径 & OS 决策点 \\
\midrule
定时器到期 & local APIC timer、HPET 或平台定时器达到比较值，产生中断请求 & 当前 CPU 是否允许接收中断，是否处在不可抢占区 \\
中断投递 & 中断控制器把 vector 投递给目标 CPU & vector 只说明入口类型，不说明调度结果 \\
入口保存 & CPU 保存返回位置、状态和部分现场，切换到内核入口和内核栈 & trap frame 必须足够恢复被打断代码 \\
内核计时 & handler 更新 tick、hrtimer、运行时间统计和超时队列 & 哪些等待者到期，当前线程时间片是否耗尽 \\
调度判断 & 若允许抢占，scheduler 比较 runqueue，选择下一个 task & 不能在持锁、关中断或不可抢占区域随意切换 \\
地址空间切换 & 若换到另一个进程，切换页表根或地址空间标识，处理必要的 TLB 状态 & 旧进程的用户虚拟地址不能被新进程继续解释 \\
执行现场切换 & 保存旧线程 callee-saved 寄存器、栈指针等，恢复新线程现场 & 漏保存一个寄存器就会破坏用户程序语义 \\
返回执行 & 中断返回指令恢复目标 task 的 RIP、状态和特权级 & 返回点可能是原线程，也可能是另一个线程上次停下的位置 \\
\bottomrule
\end{longtable}

这一案例把 CPU、定时器、中断控制器、内核栈、调度器、页表、TLB 和 cache 热状态放到同一条线上。切换页表会影响 TLB 命中；切到另一个 CPU 或另一个工作集，会影响 cache 命中；但这些只是成本的一部分。更根本的不变量是：被打断线程的现场必须完整保存，新线程的地址空间必须正确生效，中断控制器必须收到合适的确认，内核不能在不允许抢占的硬件和软件边界上强行切换。

若定时器路径出错，症状也各不相同。定时器未重新编程，超时和抢占不发生；中断被屏蔽，CPU 长时间不进入调度；EOI 顺序错误，可能出现中断风暴或后续中断被卡住；现场保存错误，会变成随机寄存器损坏；页表切换错误，则会让一个进程用另一个进程的地址解释 load/store。把这些都归结为“调度有 bug”，无法定位真正硬件边界。

\medskip
\noindent\textbf{第九条完整路径：\code{fork} 后第一次写 COW 页为什么能被精确拦住。}

\code{fork} 后，父子进程常常先共享同一批物理页。为了避免立刻复制所有内存，内核把这些页同时映射到父子进程，但清掉 writable 位。等某一方第一次写入时，硬件保护异常把这个动作拦住，内核再复制页面。这就是 copy-on-write，简称 COW。

\begin{lstlisting}
after fork:
  parent PTE: PA=X, read-only, COW
  child  PTE: PA=X, read-only, COW

child writes *p:
  hardware sees write to read-only PTE
  page fault enters kernel before store commits
\end{lstlisting}

\begin{longtable}{p{0.16\linewidth}p{0.43\linewidth}p{0.29\linewidth}}
\toprule
阶段 & 硬件和内核动作 & 为什么不能省略 \\
\midrule
建立共享 & 父子页表都指向同一物理页，PTE 清 writable，并记录 COW 语义 & 只读保护是让硬件拦住第一次写的关键 \\
清旧翻译 & 对已运行过的地址空间失效旧 writable TLB entry & 否则 CPU 可能继续用旧可写权限 \\
用户写入 & store 指令地址翻译时发现 PTE 不允许写 & 错误 store 还未提交，原指令可重试 \\
进入缺页 & CPU 保存 fault address、error code、RIP，进入内核缺页入口 & 内核能知道是哪一个地址、哪一种权限失败 \\
分配新页 & 内核分配物理页，把旧页内容复制过去，更新引用计数 & 父子后续写入不能互相污染 \\
更新 PTE & 当前进程 PTE 改成新物理页并恢复 writable，必要时 flush TLB & 新翻译必须替换旧只读翻译 \\
重试 store & 缺页处理返回后，CPU 重新执行原 store，这次写入新页 & 用户态看见的是一次普通写入，没有看到中间复制 \\
\bottomrule
\end{longtable}

这个案例容易被讲成“页表 + 缺页”两个词，但完整链路更长。写指令先经过地址生成和权限检查，TLB 可能命中旧只读翻译，也可能 page walk 读到只读 PTE；无论哪条路径，只要权限拒绝，CPU 都要在提交前进入异常。内核进入后还要确认 VMA 是否允许写、PTE 是否真是 COW、物理页引用计数是否大于一。若不是 COW，而是用户写只读代码页，就不能复制后放行，而要给进程发段错误。

COW 也说明 cache 不是唯一变量。复制旧页时，CPU 会把旧页 cache line 读入并写新页；更新 PTE 后，要处理 TLB 旧翻译；多核下，其他 CPU 可能也运行同一地址空间，需要 shootdown；内存分配器要保证新物理页不含别的进程残留数据；异常返回要保证原 store 可重试。少看任何一个硬件块，这个“懒复制”就解释不完整。

\medskip
\noindent\textbf{第十条完整路径：网卡收到一个包怎样走到 \code{recv}。}

最后看外部输入路径。网卡收到一个以太网帧，用户进程稍后调用：

\begin{lstlisting}
recv(sock, buf, len);
\end{lstlisting}

这个路径和 NVMe 读类似，都有 DMA、IOMMU、completion 和中断；但它还多了网络协议栈、socket 接收队列和从内核 buffer 到用户 buffer 的复制。它能说明“设备主动把数据带进系统”时，硬件和 OS 怎样共同维护所有权。

\begin{longtable}{p{0.15\linewidth}p{0.44\linewidth}p{0.29\linewidth}}
\toprule
阶段 & 硬件或内核动作 & 关键边界 \\
\midrule
预先布置 RX ring & 驱动分配 packet buffer，建立 DMA 映射，把 IOVA 写进接收描述符 & 设备只能 DMA 到这些已授权 buffer \\
包到达网卡 & 网卡 PHY/MAC 接收帧，控制器选择一个可用 RX descriptor & 没有空 descriptor 时可能丢包或触发流控 \\
DMA 写入内存 & 网卡经 PCIe/root complex/IOMMU 把包数据写入 buffer，并更新 descriptor 状态 & IOMMU 限制设备不能写到任意内存 \\
投递中断或轮询 & 网卡发 MSI-X，或 NAPI 轮询稍后批量收包 & 中断只是提示有工作，结果仍要读 ring 状态 \\
驱动收割 completion & CPU 读取 RX ring，做必要的 cache 同步，构造 skb 或等价包对象 & buffer 所有权从设备转回内核 \\
协议栈处理 & 内核解析二层、IP、TCP/UDP，把数据挂到 socket 接收队列 & 包可能被丢弃、重组、校验失败或排队等待 \\
用户 \code{recv} & 系统调用进入内核，检查用户 \code{buf}，把 socket 数据复制到用户页 & 用户地址可能缺页或权限失败 \\
补充 buffer & 驱动给 RX ring 重新挂上空 buffer，并写 doorbell 或更新 tail & 不补 buffer，设备后续没有位置接收新包 \\
\bottomrule
\end{longtable}

这条路径同时存在三类“内存”。第一类是普通内核对象，例如 socket 队列、skb 元数据和协议栈状态，它们由 CPU 读写并受 cache 一致性影响。第二类是 DMA buffer，提交给设备后由设备写入，completion 前 CPU 不能把它当普通可释放内存。第三类是用户 \code{buf}，它属于用户地址空间，只有系统调用检查和复制阶段才被内核安全接触。把这三类地址混在一起，就会出现把用户指针塞给设备、completion 前释放 DMA buffer、或协议栈读到未同步数据这类严重错误。

网卡路径也能说明为什么高性能 I/O 经常选择轮询和批量。每个包都发一次中断，CPU 会在中断入口、cache 扰动和协议栈调度上付出高成本；NAPI 这类机制会在中断提示后暂时转为轮询，一次收割多个 completion。优化并没有取消硬件边界：RX ring 仍要有所有权位，IOMMU 映射仍要有效，MSI-X 或轮询仍要保证内核最终发现 completion，用户 \code{recv} 仍要通过页表和权限检查。

\section{五、回到输出一个字节}

回到开头那两行：

\begin{lstlisting}
output_one_byte('A'):
  wait until the output device can accept one byte
  place 'A' into the device's input slot
\end{lstlisting}

经过前面的推导，这个需求可以对应到完整硬件路径：

\begin{longtable}{p{0.2\linewidth}p{0.38\linewidth}p{0.3\linewidth}}
\toprule
需求片段 & 硬件含义 & 对后续 OS 章节的连接 \\
\midrule
当前执行到哪一步 & CPU 从程序计数器指向的指令继续执行，参数和临时值在寄存器或栈里 & task、trap frame、上下文切换 \\
设备是否可接收 & CPU 对设备状态地址发起 load，设备状态机返回 ready/busy & 驱动轮询、超时、错误证据 \\
根据状态选择动作 & 条件分支依赖 ALU 标志位和程序计数器更新 & 条件分支、异常边界、可恢复现场 \\
把字节交给设备 & CPU 对设备寄存器对应地址发起 store，写入顺序就是命令顺序 & MMIO 顺序、屏障、posted write \\
设备自己完成输出 & 设备按自己的时钟发送字节，CPU 不参与每个 bit & 异步设备、等待队列、调度让出 CPU \\
完成或失败怎样回来 & 设备更新状态位，复杂设备还会写 completion 或发中断 & IRQ、completion、request 生命周期 \\
输出一大块数据 & CPU 准备描述符和 DMA buffer，让设备批量搬运 & DMA mapping、IOMMU、buffer ownership \\
\bottomrule
\end{longtable}

最后把它按一次真实执行走完。假设内核已经把设备 BAR 映射到内核虚拟地址 \code{dev->mmio}，其中 \code{STATUS} 偏移为 \code{0x00}，\code{DATA} 偏移为 \code{0x04}。当前 CPU 的 \code{RIP} 指向驱动里的轮询代码，某个寄存器保存 \code{dev->mmio}，另一个寄存器保存待写字节 \code{'A'}。硬件不会看到“输出一个字符”这个高级意图，它只会执行一串 load、test、branch 和 store。

\begin{longtable}{p{0.15\linewidth}p{0.42\linewidth}p{0.31\linewidth}}
\toprule
步骤 & 硬件结构路径 & 这一刻的关键边界 \\
\midrule
取指 & PC/RIP 进入取指前端，指令 cache 返回驱动代码字节 & CPU 知道下一条机器指令，不知道 C 函数语义 \\
算状态地址 & 寄存器文件读出 \code{dev->mmio}，ALU/地址生成单元加上 \code{STATUS} 偏移 & 地址只是候选，还要经过翻译和属性检查 \\
翻译地址 & MMU/TLB 把内核虚拟地址翻译成设备 BAR 的物理地址，并确认内核可访问 & 页表权限和内存类型决定这是设备访问，不是普通 RAM \\
读状态 & 总线把 MMIO read 路由到设备，设备状态寄存器返回 \code{READY/ERR} 位 & 读到的是设备保存的当前状态，不能用普通变量缓存替代 \\
条件分支 & ALU/test 逻辑检查 \code{READY}，下一 PC 选择循环或写数据路径 & 分支依赖状态位，错误时必须能保留可恢复现场 \\
算数据地址 & 地址生成单元计算 \code{dev->mmio + DATA} & 仍是 MMIO 地址，写入会变成设备命令 \\
写数据 & CPU 发出 MMIO write，设备数据寄存器接收 \code{'A'}，设备清 \code{READY} 开始工作 & store 提交后，命令已经越过 CPU/设备边界 \\
设备推进 & 设备按自己的状态机发送字节，完成后置 \code{READY} 或 \code{ERR}，复杂设备还会 IRQ & CPU 不控制每个外部 bit，只能观察状态或事件 \\
\bottomrule
\end{longtable}

这条路径里任何一步都可能改变 OS 的处理方式。取指或执行阶段若遇到异常，内核得到的是 CPU 异常现场；地址翻译失败，说明驱动映射或权限有问题；MMIO read 长期看不到 ready，驱动要按超时处理；MMIO write 后设备报错，驱动要清状态、reset 或向上层返回失败；如果改成大块输出，单个 \code{DATA} 寄存器就不够，必须升级到 descriptor、DMA、completion 和 IOMMU。也就是说，OS 不是把一个“输出函数”塞给硬件，而是在每一个硬件边界上维护状态、权限、顺序和失败证据。

结论可以概括为一条检查法：分析任意 OS 机制时，都要说明它来自哪一个机器需求，又落在哪个硬件事实上。进程能被暂停，是因为执行现场能保存；虚拟内存能懒分配，是因为缺页异常可重试；驱动要处理 timeout，是因为设备是状态机；DMA 需要严格约束，是因为设备能越过 CPU 直接访问内存；启动阶段特殊，是因为这些契约还没完全建立。

硬件链路可以收束成一张回查地图。遇到任意 OS 机制时，都可以把它放回这些机器事实：它依赖哪种状态，哪条数据通路，哪种权限检查，哪类异步事件，哪种完成证据。

\begin{longtable}{p{0.16\linewidth}p{0.38\linewidth}p{0.34\linewidth}}
\toprule
硬件事实 & 解决的问题 & 后续衔接位置 \\
\midrule
时钟/复位 & 状态寄存器让机器有可解释的上一刻和下一刻 & 启动、异常返回、设备 reset 和错误恢复 \\
提交边界 & ALU、MUX 和寄存器让指令结果从候选值变成软件可见状态 & trap frame、上下文切换、精确异常和调试证据 \\
入口现场 & 程序计数器、栈和 ABI 保存普通调用和异常入口的返回路径 & 系统调用、信号、线程切换和 \code{exec} 用户栈 \\
地址翻译 & 页表、TLB、cache 和内存属性让地址访问能被翻译、缓存、拒绝或重试 & 虚拟内存、page fault、COW、TLB shootdown 和锁性能 \\
设备寄存器 & MMIO、总线、设备状态位和中断让外部设备进入 CPU 世界 & 驱动、定时器、等待队列、超时和中断下半部 \\
DMA 完成 & descriptor、DMA、IOMMU 和 completion 让设备批量搬运数据，并把完成事件对应回请求 & 块层、网卡、page cache、buffer 生命周期和错误返回 \\
多核同步 & 一致性协议和原子操作让多个执行流共享内存但仍能维护不变量 & 调度、锁、futex、RCU、伪共享和内存屏障 \\
平台资源 & 固件、ACPI/PCIe、NVMe/FTL 和安全硬件让内核发现资源并处理故障证据 & 设备枚举、持久化、崩溃恢复、权限隔离和观测 \\
\bottomrule
\end{longtable}

遇到一个 OS 名词，不应只背定义，还要追问它压在哪个硬件事实上。进程压在可保存执行现场上；虚拟内存压在页表、TLB 和 fault 上；文件 I/O 压在 page cache、descriptor、DMA 和 completion 上；安全隔离压在特权级、页表权限、对象引用和审计证据上。这样硬件基础不会变成前置名词表，而会持续参与每一个机制解释。

第二章到这里结束。后面不再继续堆硬件名词，而是换一个视角：既然硬件只提供现场、地址、权限、队列、事件和完成证据，内核为什么必须把它们包装成 task、地址空间、文件对象、request、wait queue 和 completion。下一章就从一次写日志请求开始，把这些硬件事实收束成内核对象。
